%%%%%%%%%%%%%%%%%%%%%%%%%%%%%%%%%%%%%%%%%%%%%%%%%%%%%%%%%%%%
% 13.9 MATHEMATICAL ECONOMICS
% The Composition of Advanced Mathematics
%%%%%%%%%%%%%%%%%%%%%%%%%%%%%%%%%%%%%%%%%%%%%%%%%%%%%%%%%%%%

\section{Mathematical Economics}
\label{sec:mathematical-economics}

\prerequisite{
Calculus, multivariable calculus, linear algebra, optimization, probability,
statistics, differential equations, game theory, mathematical modeling,
and mathematical finance.
}

\learningobjectives{
By the end of this section, the reader should be able to:
\begin{itemize}
    \item explain the role of mathematics in economics;
    \item distinguish positive and normative economic models;
    \item formulate economic variables, parameters, and constraints;
    \item understand marginal analysis;
    \item compute marginal cost, revenue, and utility;
    \item understand elasticity;
    \item compute price elasticity of demand;
    \item distinguish elastic, inelastic, and unit-elastic demand;
    \item understand consumer preference relations;
    \item define utility functions;
    \item recognize ordinal utility;
    \item understand indifference curves;
    \item formulate utility-maximization problems;
    \item derive first-order conditions for consumer choice;
    \item understand marginal utility;
    \item define marginal rate of substitution;
    \item understand budget constraints;
    \item derive Marshallian demand conceptually;
    \item understand indirect utility conceptually;
    \item recognize expenditure minimization;
    \item understand Hicksian demand conceptually;
    \item recognize income and substitution effects;
    \item understand Slutsky decomposition conceptually;
    \item define production functions;
    \item recognize Cobb--Douglas production;
    \item understand marginal products;
    \item understand returns to scale;
    \item formulate cost-minimization problems;
    \item understand isoquants;
    \item define marginal rate of technical substitution;
    \item derive conditional factor demand conceptually;
    \item formulate profit-maximization problems;
    \item understand competitive-firm supply;
    \item distinguish fixed and variable costs;
    \item understand average and marginal cost;
    \item recognize economies and diseconomies of scale;
    \item formulate market demand and supply;
    \item solve competitive market equilibrium;
    \item analyze comparative statics;
    \item understand consumer and producer surplus;
    \item recognize welfare measures;
    \item understand taxes and subsidies mathematically;
    \item recognize price ceilings and price floors;
    \item understand monopoly optimization;
    \item derive marginal-revenue conditions;
    \item understand price discrimination conceptually;
    \item recognize oligopoly models;
    \item understand Cournot competition;
    \item understand Bertrand competition conceptually;
    \item recognize Stackelberg competition conceptually;
    \item understand Nash equilibrium;
    \item recognize dominant strategies;
    \item understand mixed strategies conceptually;
    \item recognize repeated games conceptually;
    \item understand public goods;
    \item recognize externalities;
    \item understand Pigouvian taxes conceptually;
    \item understand common-resource problems;
    \item recognize mechanism design conceptually;
    \item understand social choice conceptually;
    \item recognize auctions;
    \item understand expected utility;
    \item recognize risk aversion;
    \item understand certainty equivalents;
    \item define risk premiums;
    \item recognize Arrow--Pratt measures conceptually;
    \item understand intertemporal choice;
    \item recognize discounting;
    \item formulate consumption-saving problems;
    \item understand dynamic optimization;
    \item recognize Bellman equations in economics;
    \item understand growth models;
    \item recognize the Solow model;
    \item understand steady states;
    \item analyze local stability of economic equilibria;
    \item recognize overlapping-generations models conceptually;
    \item understand optimal growth conceptually;
    \item recognize differential and difference equation models in economics;
    \item understand input-output models;
    \item recognize Leontief systems;
    \item understand general equilibrium conceptually;
    \item recognize excess-demand functions;
    \item understand Walras' law conceptually;
    \item recognize fixed-point ideas in equilibrium theory;
    \item understand econometric models conceptually;
    \item distinguish structural and reduced-form equations;
    \item recognize identification problems conceptually;
    \item understand constrained optimization in economic policy;
    \item recognize network effects in economics;
    \item understand financial-economic connections;
    \item distinguish model equilibrium from real-world prediction;
    \item understand sensitivity, calibration, and uncertainty;
    \item connect mathematical economics with optimization, game theory,
          probability, statistics, dynamical systems, finance,
          and network science.
\end{itemize}
}

%%%%%%%%%%%%%%%%%%%%%%%%%%%%%%%%%%%%%%%%%%%%%%%%%%%%%%%%%%%%
\subsection{What Is Mathematical Economics?}
\label{subsec:what-is-mathematical-economics}
%%%%%%%%%%%%%%%%%%%%%%%%%%%%%%%%%%%%%%%%%%%%%%%%%%%%%%%%%%%%

Mathematical economics uses mathematical structures to describe economic
behavior and economic systems.

A model may study:

\[
\boxed{
\text{consumer choice},
}
\]

\[
\boxed{
\text{firm behavior},
}
\]

\[
\boxed{
\text{markets},
}
\]

\[
\boxed{
\text{strategic interaction},
}
\]

\[
\boxed{
\text{growth},
}
\]

or

\[
\boxed{
\text{policy}.
}
\]

The goal is not merely to calculate quantities, but to understand how
assumptions imply economic behavior.

%%%%%%%%%%%%%%%%%%%%%%%%%%%%%%%%%%%%%%%%%%%%%%%%%%%%%%%%%%%%
\subsection{Economic Variables and Parameters}
\label{subsec:economic-variables-parameters}
%%%%%%%%%%%%%%%%%%%%%%%%%%%%%%%%%%%%%%%%%%%%%%%%%%%%%%%%%%%%

An economic model distinguishes endogenous variables from exogenous
parameters.

For example,

\[
\boxed{
Q_D
=
Q_D(P,Y)
}
\]

may describe quantity demanded as a function of price \(P\) and income \(Y\).

Here \(Q_D\) is endogenous once \(P\) and \(Y\) are specified.

%%%%%%%%%%%%%%%%%%%%%%%%%%%%%%%%%%%%%%%%%%%%%%%%%%%%%%%%%%%%
\subsection{Marginal Analysis}
\label{subsec:marginal-analysis}
%%%%%%%%%%%%%%%%%%%%%%%%%%%%%%%%%%%%%%%%%%%%%%%%%%%%%%%%%%%%

The word marginal refers to a derivative or incremental change.

If total cost is

\[
C(q),
\]

then marginal cost is

\[
\boxed{
MC(q)
=
C'(q).
}
\]

If revenue is

\[
R(q),
\]

then marginal revenue is

\[
\boxed{
MR(q)
=
R'(q).
}
\]

%%%%%%%%%%%%%%%%%%%%%%%%%%%%%%%%%%%%%%%%%%%%%%%%%%%%%%%%%%%%
\subsection{Profit}
\label{subsec:economic-profit}
%%%%%%%%%%%%%%%%%%%%%%%%%%%%%%%%%%%%%%%%%%%%%%%%%%%%%%%%%%%%

Profit is

\[
\boxed{
\pi(q)
=
R(q)-C(q).
}
\]

An interior profit maximum satisfies

\[
\boxed{
\pi'(q)=0.
}
\]

Thus

\[
R'(q)-C'(q)=0,
\]

or

\[
\boxed{
MR(q)=MC(q).
}
\]

%%%%%%%%%%%%%%%%%%%%%%%%%%%%%%%%%%%%%%%%%%%%%%%%%%%%%%%%%%%%
\subsection{Worked Example: Profit Maximization}
\label{subsec:worked-profit-maximization}
%%%%%%%%%%%%%%%%%%%%%%%%%%%%%%%%%%%%%%%%%%%%%%%%%%%%%%%%%%%%

\begin{workedexamplebox}{A Simple Firm}

Suppose

\[
R(q)=40q-q^2
\]

and

\[
C(q)=10q+50.
\]

Then

\[
\pi(q)
=
30q-q^2-50.
\]

Differentiate:

\[
\pi'(q)
=
30-2q.
\]

Set

\[
30-2q=0.
\]

Thus

\[
q^*=15.
\]

Since

\[
\pi''(q)=-2<0,
\]

this is a maximum.

\begin{answerbox}
\[
\boxed{
q^*=15.
}
\]
\end{answerbox}

\end{workedexamplebox}

%%%%%%%%%%%%%%%%%%%%%%%%%%%%%%%%%%%%%%%%%%%%%%%%%%%%%%%%%%%%
\subsection{Elasticity}
\label{subsec:economic-elasticity}
%%%%%%%%%%%%%%%%%%%%%%%%%%%%%%%%%%%%%%%%%%%%%%%%%%%%%%%%%%%%

Elasticity measures proportional responsiveness.

If

\[
y=f(x),
\]

the elasticity of \(y\) with respect to \(x\) is

\[
\boxed{
\varepsilon_{y,x}
=
\frac{dy}{dx}
\frac{x}{y}.
}
\]

It compares percentage changes rather than absolute changes.

%%%%%%%%%%%%%%%%%%%%%%%%%%%%%%%%%%%%%%%%%%%%%%%%%%%%%%%%%%%%
\subsection{Price Elasticity of Demand}
\label{subsec:price-elasticity-demand}
%%%%%%%%%%%%%%%%%%%%%%%%%%%%%%%%%%%%%%%%%%%%%%%%%%%%%%%%%%%%

If quantity demanded is

\[
Q=Q(P),
\]

then price elasticity is

\[
\boxed{
\varepsilon_D
=
\frac{dQ}{dP}
\frac{P}{Q}.
}
\]

Because demand often decreases with price,

\[
\varepsilon_D<0.
\]

Many texts report its absolute magnitude.

%%%%%%%%%%%%%%%%%%%%%%%%%%%%%%%%%%%%%%%%%%%%%%%%%%%%%%%%%%%%
\subsection{Elastic, Inelastic, and Unit Elastic}
\label{subsec:elasticity-classification}
%%%%%%%%%%%%%%%%%%%%%%%%%%%%%%%%%%%%%%%%%%%%%%%%%%%%%%%%%%%%

Using absolute elasticity,

\[
\boxed{
|\varepsilon_D|>1
}
\]

means elastic demand,

\[
\boxed{
|\varepsilon_D|<1
}
\]

means inelastic demand, and

\[
\boxed{
|\varepsilon_D|=1
}
\]

means unit-elastic demand.

%%%%%%%%%%%%%%%%%%%%%%%%%%%%%%%%%%%%%%%%%%%%%%%%%%%%%%%%%%%%
\subsection{Worked Example: Demand Elasticity}
\label{subsec:worked-demand-elasticity}
%%%%%%%%%%%%%%%%%%%%%%%%%%%%%%%%%%%%%%%%%%%%%%%%%%%%%%%%%%%%

\begin{workedexamplebox}{Linear Demand}

Suppose

\[
Q(P)=100-2P.
\]

Then

\[
\frac{dQ}{dP}
=
-2.
\]

At

\[
P=20,
\]

we have

\[
Q=60.
\]

Therefore,

\[
\varepsilon_D
=
(-2)\frac{20}{60}
=
-\frac23.
\]

\begin{answerbox}
\[
\boxed{
\varepsilon_D=-\frac23,
}
\]

so demand is inelastic at this point.
\end{answerbox}

\end{workedexamplebox}

%%%%%%%%%%%%%%%%%%%%%%%%%%%%%%%%%%%%%%%%%%%%%%%%%%%%%%%%%%%%
\subsection{Preferences}
\label{subsec:consumer-preferences}
%%%%%%%%%%%%%%%%%%%%%%%%%%%%%%%%%%%%%%%%%%%%%%%%%%%%%%%%%%%%

Consumer theory begins with a preference relation over bundles.

Write

\[
\boxed{
\mathbf{x}\succeq\mathbf{y}
}
\]

to mean bundle \(\mathbf{x}\) is at least as preferred as bundle
\(\mathbf{y}\).

Common assumptions include:

\[
\boxed{
\text{completeness}
}
\]

and

\[
\boxed{
\text{transitivity}.
}
\]

%%%%%%%%%%%%%%%%%%%%%%%%%%%%%%%%%%%%%%%%%%%%%%%%%%%%%%%%%%%%
\subsection{Utility Functions}
\label{subsec:utility-functions}
%%%%%%%%%%%%%%%%%%%%%%%%%%%%%%%%%%%%%%%%%%%%%%%%%%%%%%%%%%%%

A utility function assigns numbers to bundles while preserving preference
ordering.

\[
\boxed{
u(\mathbf{x})
\geq
u(\mathbf{y})
}
\]

if and only if

\[
\mathbf{x}
\succeq
\mathbf{y}
\]

under an appropriate representation.

Utility is ordinarily interpreted ordinally rather than as a direct physical
measurement of satisfaction.

%%%%%%%%%%%%%%%%%%%%%%%%%%%%%%%%%%%%%%%%%%%%%%%%%%%%%%%%%%%%
\subsection{Indifference Curves}
\label{subsec:indifference-curves}
%%%%%%%%%%%%%%%%%%%%%%%%%%%%%%%%%%%%%%%%%%%%%%%%%%%%%%%%%%%%

An indifference curve is a level set of utility:

\[
\boxed{
u(x_1,x_2)=\bar u.
}
\]

Every bundle on the curve yields the same represented preference level.

Its slope reflects the tradeoff between goods.

%%%%%%%%%%%%%%%%%%%%%%%%%%%%%%%%%%%%%%%%%%%%%%%%%%%%%%%%%%%%
\subsection{Marginal Utility}
\label{subsec:marginal-utility}
%%%%%%%%%%%%%%%%%%%%%%%%%%%%%%%%%%%%%%%%%%%%%%%%%%%%%%%%%%%%

For utility

\[
u(x_1,\ldots,x_n),
\]

marginal utility of good \(i\) is

\[
\boxed{
MU_i
=
\frac{\partial u}{\partial x_i}.
}
\]

It measures local utility change from increasing consumption of good \(i\),
holding other goods fixed.

%%%%%%%%%%%%%%%%%%%%%%%%%%%%%%%%%%%%%%%%%%%%%%%%%%%%%%%%%%%%
\subsection{Marginal Rate of Substitution}
\label{subsec:marginal-rate-substitution}
%%%%%%%%%%%%%%%%%%%%%%%%%%%%%%%%%%%%%%%%%%%%%%%%%%%%%%%%%%%%

For two goods, along an indifference curve,

\[
du=0.
\]

Thus

\[
u_{x_1}\,dx_1
+
u_{x_2}\,dx_2
=
0.
\]

Therefore,

\[
\boxed{
-\frac{dx_2}{dx_1}
=
\frac{MU_1}{MU_2}.
}
\]

This ratio is the marginal rate of substitution.

%%%%%%%%%%%%%%%%%%%%%%%%%%%%%%%%%%%%%%%%%%%%%%%%%%%%%%%%%%%%
\subsection{Budget Constraint}
\label{subsec:consumer-budget-constraint}
%%%%%%%%%%%%%%%%%%%%%%%%%%%%%%%%%%%%%%%%%%%%%%%%%%%%%%%%%%%%

For prices \(p_1,p_2\) and income \(m\),

\[
\boxed{
p_1x_1+p_2x_2
\leq
m.
}
\]

For monotonic preferences, an interior optimum typically uses the full
budget:

\[
\boxed{
p_1x_1+p_2x_2=m.
}
\]

%%%%%%%%%%%%%%%%%%%%%%%%%%%%%%%%%%%%%%%%%%%%%%%%%%%%%%%%%%%%
\subsection{Utility Maximization}
\label{subsec:utility-maximization}
%%%%%%%%%%%%%%%%%%%%%%%%%%%%%%%%%%%%%%%%%%%%%%%%%%%%%%%%%%%%

The consumer problem is

\[
\boxed{
\max_{\mathbf{x}}
\quad
u(\mathbf{x})
}
\]

subject to

\[
\boxed{
\mathbf{p}^T\mathbf{x}
\leq
m,
\qquad
\mathbf{x}\geq0.
}
\]

For an interior solution, construct

\[
\boxed{
\mathcal{L}
=
u(\mathbf{x})
+
\lambda
\left(
m-\mathbf{p}^T\mathbf{x}
\right).
}
\]

%%%%%%%%%%%%%%%%%%%%%%%%%%%%%%%%%%%%%%%%%%%%%%%%%%%%%%%%%%%%
\subsection{Consumer First-Order Conditions}
\label{subsec:consumer-first-order-conditions}
%%%%%%%%%%%%%%%%%%%%%%%%%%%%%%%%%%%%%%%%%%%%%%%%%%%%%%%%%%%%

For an interior optimum,

\[
\boxed{
\frac{\partial u}{\partial x_i}
=
\lambda p_i.
}
\]

For two goods,

\[
\frac{MU_1}{p_1}
=
\frac{MU_2}{p_2}.
\]

Equivalently,

\[
\boxed{
\frac{MU_1}{MU_2}
=
\frac{p_1}{p_2}.
}
\]

The marginal rate of substitution equals the price ratio.

%%%%%%%%%%%%%%%%%%%%%%%%%%%%%%%%%%%%%%%%%%%%%%%%%%%%%%%%%%%%
\subsection{Cobb--Douglas Utility}
\label{subsec:cobb-douglas-utility}
%%%%%%%%%%%%%%%%%%%%%%%%%%%%%%%%%%%%%%%%%%%%%%%%%%%%%%%%%%%%

A common utility function is

\[
\boxed{
u(x,y)
=
x^\alpha y^\beta,
}
\]

where

\[
\alpha,\beta>0.
\]

Often one uses

\[
\alpha+\beta=1,
\]

although this normalization is not necessary for the preference ordering.

%%%%%%%%%%%%%%%%%%%%%%%%%%%%%%%%%%%%%%%%%%%%%%%%%%%%%%%%%%%%
\subsection{Worked Example: Cobb--Douglas Demand}
\label{subsec:worked-cobb-douglas-demand}
%%%%%%%%%%%%%%%%%%%%%%%%%%%%%%%%%%%%%%%%%%%%%%%%%%%%%%%%%%%%

\begin{workedexamplebox}{Utility Maximization}

Suppose

\[
u(x,y)=x^{1/2}y^{1/2}
\]

and

\[
p_xx+p_yy=m.
\]

For this symmetric utility, the optimal expenditure shares are equal.

Thus

\[
p_xx
=
\frac{m}{2},
\]

and

\[
p_yy
=
\frac{m}{2}.
\]

Therefore,

\begin{answerbox}
\[
\boxed{
x^*
=
\frac{m}{2p_x},
\qquad
y^*
=
\frac{m}{2p_y}.
}
\]
\end{answerbox}

\end{workedexamplebox}

%%%%%%%%%%%%%%%%%%%%%%%%%%%%%%%%%%%%%%%%%%%%%%%%%%%%%%%%%%%%
\subsection{Marshallian Demand}
\label{subsec:marshallian-demand}
%%%%%%%%%%%%%%%%%%%%%%%%%%%%%%%%%%%%%%%%%%%%%%%%%%%%%%%%%%%%

The utility-maximizing choice is called Marshallian or uncompensated demand:

\[
\boxed{
\mathbf{x}
=
\mathbf{x}(\mathbf{p},m).
}
\]

It depends on prices and income.

Marshallian demand incorporates both substitution and income effects.

%%%%%%%%%%%%%%%%%%%%%%%%%%%%%%%%%%%%%%%%%%%%%%%%%%%%%%%%%%%%
\subsection{Indirect Utility}
\label{subsec:indirect-utility}
%%%%%%%%%%%%%%%%%%%%%%%%%%%%%%%%%%%%%%%%%%%%%%%%%%%%%%%%%%%%

The indirect utility function gives maximum attainable utility at prices
\(\mathbf{p}\) and income \(m\):

\[
\boxed{
v(\mathbf{p},m)
=
\max_{\mathbf{x}}
\left\{
u(\mathbf{x})
:
\mathbf{p}^T\mathbf{x}\leq m
\right\}.
}
\]

It converts the consumer optimization problem into a function of market
conditions.

%%%%%%%%%%%%%%%%%%%%%%%%%%%%%%%%%%%%%%%%%%%%%%%%%%%%%%%%%%%%
\subsection{Expenditure Minimization}
\label{subsec:expenditure-minimization}
%%%%%%%%%%%%%%%%%%%%%%%%%%%%%%%%%%%%%%%%%%%%%%%%%%%%%%%%%%%%

The dual consumer problem is

\[
\boxed{
\min_{\mathbf{x}}
\quad
\mathbf{p}^T\mathbf{x}
}
\]

subject to

\[
\boxed{
u(\mathbf{x})
\geq
\bar u.
}
\]

The minimum expenditure is the expenditure function

\[
\boxed{
e(\mathbf{p},\bar u).
}
\]

%%%%%%%%%%%%%%%%%%%%%%%%%%%%%%%%%%%%%%%%%%%%%%%%%%%%%%%%%%%%
\subsection{Hicksian Demand}
\label{subsec:hicksian-demand}
%%%%%%%%%%%%%%%%%%%%%%%%%%%%%%%%%%%%%%%%%%%%%%%%%%%%%%%%%%%%

The expenditure-minimizing bundle is Hicksian or compensated demand:

\[
\boxed{
\mathbf{h}
=
\mathbf{h}(\mathbf{p},\bar u).
}
\]

It answers:

\[
\boxed{
\text{How much of each good minimizes spending while preserving utility?}
}
\]

%%%%%%%%%%%%%%%%%%%%%%%%%%%%%%%%%%%%%%%%%%%%%%%%%%%%%%%%%%%%
\subsection{Income and Substitution Effects}
\label{subsec:income-substitution-effects}
%%%%%%%%%%%%%%%%%%%%%%%%%%%%%%%%%%%%%%%%%%%%%%%%%%%%%%%%%%%%

A price change affects demand through two mechanisms.

The substitution effect reflects changes in relative prices while holding
utility appropriately compensated.

The income effect reflects the change in real purchasing power.

Thus

\[
\boxed{
\text{total price effect}
=
\text{substitution effect}
+
\text{income effect}.
}
\]

%%%%%%%%%%%%%%%%%%%%%%%%%%%%%%%%%%%%%%%%%%%%%%%%%%%%%%%%%%%%
\subsection{Production Functions}
\label{subsec:production-functions}
%%%%%%%%%%%%%%%%%%%%%%%%%%%%%%%%%%%%%%%%%%%%%%%%%%%%%%%%%%%%

A production function maps inputs to output:

\[
\boxed{
q
=
f(x_1,\ldots,x_n).
}
\]

Common inputs include:

\[
\boxed{
\text{labor},
}
\]

\[
\boxed{
\text{capital},
}
\]

\[
\boxed{
\text{energy},
}
\]

and

\[
\boxed{
\text{materials}.
}
\]

%%%%%%%%%%%%%%%%%%%%%%%%%%%%%%%%%%%%%%%%%%%%%%%%%%%%%%%%%%%%
\subsection{Marginal Product}
\label{subsec:marginal-product}
%%%%%%%%%%%%%%%%%%%%%%%%%%%%%%%%%%%%%%%%%%%%%%%%%%%%%%%%%%%%

Marginal product of input \(i\) is

\[
\boxed{
MP_i
=
\frac{\partial f}{\partial x_i}.
}
\]

It measures the local increase in output from one additional unit of input
\(i\), holding other inputs fixed.

%%%%%%%%%%%%%%%%%%%%%%%%%%%%%%%%%%%%%%%%%%%%%%%%%%%%%%%%%%%%
\subsection{Cobb--Douglas Production}
\label{subsec:cobb-douglas-production}
%%%%%%%%%%%%%%%%%%%%%%%%%%%%%%%%%%%%%%%%%%%%%%%%%%%%%%%%%%%%

A common production function is

\[
\boxed{
q
=
AK^\alpha L^\beta,
}
\]

where

\[
A>0
\]

is productivity,

\[
K
\]

is capital, and

\[
L
\]

is labor.

%%%%%%%%%%%%%%%%%%%%%%%%%%%%%%%%%%%%%%%%%%%%%%%%%%%%%%%%%%%%
\subsection{Returns to Scale}
\label{subsec:returns-to-scale}
%%%%%%%%%%%%%%%%%%%%%%%%%%%%%%%%%%%%%%%%%%%%%%%%%%%%%%%%%%%%

Consider scaling all inputs by \(t>0\).

If

\[
f(t\mathbf{x})
=
t^rf(\mathbf{x}),
\]

then \(f\) is homogeneous of degree \(r\).

For Cobb--Douglas,

\[
\boxed{
r=\alpha+\beta.
}
\]

Therefore:

\[
\boxed{
\alpha+\beta>1
}
\]

gives increasing returns to scale,

\[
\boxed{
\alpha+\beta=1
}
\]

gives constant returns to scale,

and

\[
\boxed{
\alpha+\beta<1
}
\]

gives decreasing returns to scale.

%%%%%%%%%%%%%%%%%%%%%%%%%%%%%%%%%%%%%%%%%%%%%%%%%%%%%%%%%%%%
\subsection{Isoquants}
\label{subsec:isoquants}
%%%%%%%%%%%%%%%%%%%%%%%%%%%%%%%%%%%%%%%%%%%%%%%%%%%%%%%%%%%%

An isoquant is a level set of production:

\[
\boxed{
f(K,L)=\bar q.
}
\]

All points on an isoquant produce the same output.

Its slope reflects substitutability between inputs.

%%%%%%%%%%%%%%%%%%%%%%%%%%%%%%%%%%%%%%%%%%%%%%%%%%%%%%%%%%%%
\subsection{Marginal Rate of Technical Substitution}
\label{subsec:mrts}
%%%%%%%%%%%%%%%%%%%%%%%%%%%%%%%%%%%%%%%%%%%%%%%%%%%%%%%%%%%%

Along an isoquant,

\[
df=0.
\]

Thus

\[
f_K\,dK
+
f_L\,dL
=
0.
\]

Therefore,

\[
\boxed{
-\frac{dK}{dL}
=
\frac{MP_L}{MP_K}.
}
\]

This is a marginal rate of technical substitution.

%%%%%%%%%%%%%%%%%%%%%%%%%%%%%%%%%%%%%%%%%%%%%%%%%%%%%%%%%%%%
\subsection{Cost Minimization}
\label{subsec:firm-cost-minimization}
%%%%%%%%%%%%%%%%%%%%%%%%%%%%%%%%%%%%%%%%%%%%%%%%%%%%%%%%%%%%

Suppose capital and labor prices are \(r\) and \(w\).

The firm solves

\[
\boxed{
\min_{K,L}
\quad
rK+wL
}
\]

subject to

\[
\boxed{
f(K,L)
\geq
\bar q.
}
\]

At an interior optimum,

\[
\boxed{
\frac{MP_L}{MP_K}
=
\frac{w}{r}.
}
\]

%%%%%%%%%%%%%%%%%%%%%%%%%%%%%%%%%%%%%%%%%%%%%%%%%%%%%%%%%%%%
\subsection{Cost Function}
\label{subsec:economic-cost-function}
%%%%%%%%%%%%%%%%%%%%%%%%%%%%%%%%%%%%%%%%%%%%%%%%%%%%%%%%%%%%

The cost function gives minimum expenditure needed to produce output \(q\):

\[
\boxed{
C(\mathbf{w},q)
=
\min_{\mathbf{x}}
\left\{
\mathbf{w}^T\mathbf{x}
:
f(\mathbf{x})\geq q
\right\}.
}
\]

It summarizes technology and input prices.

%%%%%%%%%%%%%%%%%%%%%%%%%%%%%%%%%%%%%%%%%%%%%%%%%%%%%%%%%%%%
\subsection{Average and Marginal Cost}
\label{subsec:average-marginal-cost}
%%%%%%%%%%%%%%%%%%%%%%%%%%%%%%%%%%%%%%%%%%%%%%%%%%%%%%%%%%%%

Average cost is

\[
\boxed{
AC(q)
=
\frac{C(q)}{q}.
}
\]

Marginal cost is

\[
\boxed{
MC(q)
=
C'(q).
}
\]

When marginal cost lies below average cost, average cost decreases.

When marginal cost lies above average cost, average cost increases.

%%%%%%%%%%%%%%%%%%%%%%%%%%%%%%%%%%%%%%%%%%%%%%%%%%%%%%%%%%%%
\subsection{Profit Maximization for a Competitive Firm}
\label{subsec:competitive-firm-profit}
%%%%%%%%%%%%%%%%%%%%%%%%%%%%%%%%%%%%%%%%%%%%%%%%%%%%%%%%%%%%

A price-taking firm faces market price \(p\).

Profit is

\[
\boxed{
\pi(q)
=
pq-C(q).
}
\]

An interior optimum satisfies

\[
\boxed{
p=MC(q).
}
\]

Thus competitive supply is connected to the marginal-cost curve under
appropriate conditions.

%%%%%%%%%%%%%%%%%%%%%%%%%%%%%%%%%%%%%%%%%%%%%%%%%%%%%%%%%%%%
\subsection{Market Demand}
\label{subsec:market-demand}
%%%%%%%%%%%%%%%%%%%%%%%%%%%%%%%%%%%%%%%%%%%%%%%%%%%%%%%%%%%%

Market demand aggregates individual demands.

If consumer \(i\) has demand

\[
q_i(P),
\]

then

\[
\boxed{
Q_D(P)
=
\sum_iq_i(P).
}
\]

Market demand typically decreases with price under standard assumptions.

%%%%%%%%%%%%%%%%%%%%%%%%%%%%%%%%%%%%%%%%%%%%%%%%%%%%%%%%%%%%
\subsection{Market Supply}
\label{subsec:market-supply}
%%%%%%%%%%%%%%%%%%%%%%%%%%%%%%%%%%%%%%%%%%%%%%%%%%%%%%%%%%%%

Market supply aggregates firm supplies:

\[
\boxed{
Q_S(P)
=
\sum_jq_j(P).
}
\]

In many basic models, market supply increases with price.

%%%%%%%%%%%%%%%%%%%%%%%%%%%%%%%%%%%%%%%%%%%%%%%%%%%%%%%%%%%%
\subsection{Competitive Market Equilibrium}
\label{subsec:competitive-market-equilibrium}
%%%%%%%%%%%%%%%%%%%%%%%%%%%%%%%%%%%%%%%%%%%%%%%%%%%%%%%%%%%%

A competitive equilibrium price \(P^*\) satisfies

\[
\boxed{
Q_D(P^*)
=
Q_S(P^*).
}
\]

The equilibrium quantity is

\[
\boxed{
Q^*
=
Q_D(P^*)
=
Q_S(P^*).
}
\]

%%%%%%%%%%%%%%%%%%%%%%%%%%%%%%%%%%%%%%%%%%%%%%%%%%%%%%%%%%%%
\subsection{Worked Example: Market Equilibrium}
\label{subsec:worked-market-equilibrium}
%%%%%%%%%%%%%%%%%%%%%%%%%%%%%%%%%%%%%%%%%%%%%%%%%%%%%%%%%%%%

\begin{workedexamplebox}{Linear Supply and Demand}

Suppose

\[
Q_D=100-2P
\]

and

\[
Q_S=20+2P.
\]

Set

\[
100-2P
=
20+2P.
\]

Then

\[
80=4P.
\]

Therefore,

\[
P^*=20.
\]

Substituting,

\[
Q^*
=
60.
\]

\begin{answerbox}
\[
\boxed{
P^*=20,
\qquad
Q^*=60.
}
\]
\end{answerbox}

\end{workedexamplebox}

%%%%%%%%%%%%%%%%%%%%%%%%%%%%%%%%%%%%%%%%%%%%%%%%%%%%%%%%%%%%
\subsection{Comparative Statics}
\label{subsec:comparative-statics}
%%%%%%%%%%%%%%%%%%%%%%%%%%%%%%%%%%%%%%%%%%%%%%%%%%%%%%%%%%%%

Comparative statics studies how an equilibrium changes when a parameter
changes.

Suppose equilibrium satisfies

\[
\boxed{
F(x,\theta)=0.
}
\]

Under suitable differentiability and regularity conditions,

\[
\boxed{
\frac{dx^*}{d\theta}
=
-
\frac{
F_\theta
}{
F_x
}.
}
\]

This is an application of the implicit function theorem.

%%%%%%%%%%%%%%%%%%%%%%%%%%%%%%%%%%%%%%%%%%%%%%%%%%%%%%%%%%%%
\subsection{Worked Example: Comparative Statics}
\label{subsec:worked-comparative-statics}
%%%%%%%%%%%%%%%%%%%%%%%%%%%%%%%%%%%%%%%%%%%%%%%%%%%%%%%%%%%%

\begin{workedexamplebox}{Demand Shift}

Suppose

\[
Q_D=a-bP
\]

and

\[
Q_S=c+dP.
\]

Equilibrium satisfies

\[
a-bP=c+dP.
\]

Thus

\[
P^*
=
\frac{a-c}{b+d}.
\]

Differentiate with respect to \(a\):

\[
\frac{\partial P^*}{\partial a}
=
\frac{1}{b+d}.
\]

Therefore,

\begin{answerbox}
\[
\boxed{
\frac{\partial P^*}{\partial a}>0
}
\]

when \(b,d>0\).
\end{answerbox}

An outward demand shift increases equilibrium price.

\end{workedexamplebox}

%%%%%%%%%%%%%%%%%%%%%%%%%%%%%%%%%%%%%%%%%%%%%%%%%%%%%%%%%%%%
\subsection{Consumer Surplus}
\label{subsec:consumer-surplus}
%%%%%%%%%%%%%%%%%%%%%%%%%%%%%%%%%%%%%%%%%%%%%%%%%%%%%%%%%%%%

Consumer surplus measures the difference between willingness to pay and
actual expenditure.

If inverse demand is \(P_D(Q)\) and market price is \(P^*\),

\[
\boxed{
CS
=
\int_0^{Q^*}
\left[
P_D(Q)-P^*
\right]
dQ.
}
\]

%%%%%%%%%%%%%%%%%%%%%%%%%%%%%%%%%%%%%%%%%%%%%%%%%%%%%%%%%%%%
\subsection{Producer Surplus}
\label{subsec:producer-surplus}
%%%%%%%%%%%%%%%%%%%%%%%%%%%%%%%%%%%%%%%%%%%%%%%%%%%%%%%%%%%%

Producer surplus measures revenue above variable supply cost.

If inverse supply is \(P_S(Q)\),

\[
\boxed{
PS
=
\int_0^{Q^*}
\left[
P^*-P_S(Q)
\right]
dQ.
}
\]

%%%%%%%%%%%%%%%%%%%%%%%%%%%%%%%%%%%%%%%%%%%%%%%%%%%%%%%%%%%%
\subsection{Total Surplus}
\label{subsec:total-surplus}
%%%%%%%%%%%%%%%%%%%%%%%%%%%%%%%%%%%%%%%%%%%%%%%%%%%%%%%%%%%%

In a basic competitive market without externalities or other distortions,

\[
\boxed{
TS
=
CS+PS.
}
\]

This quantity is often used as a simple welfare measure.

Its interpretation depends on the assumptions of the model.

%%%%%%%%%%%%%%%%%%%%%%%%%%%%%%%%%%%%%%%%%%%%%%%%%%%%%%%%%%%%
\subsection{Taxes}
\label{subsec:economic-taxes}
%%%%%%%%%%%%%%%%%%%%%%%%%%%%%%%%%%%%%%%%%%%%%%%%%%%%%%%%%%%%

Suppose a per-unit tax \(t\) creates a wedge between buyer and seller prices:

\[
\boxed{
P_B-P_S=t.
}
\]

Equilibrium must satisfy demand at \(P_B\) and supply at \(P_S\).

The burden of the tax depends on relative elasticities rather than only on
who legally remits the tax.

%%%%%%%%%%%%%%%%%%%%%%%%%%%%%%%%%%%%%%%%%%%%%%%%%%%%%%%%%%%%
\subsection{Tax Incidence}
\label{subsec:tax-incidence}
%%%%%%%%%%%%%%%%%%%%%%%%%%%%%%%%%%%%%%%%%%%%%%%%%%%%%%%%%%%%

The side of the market that is less responsive to price generally bears a
larger share of a small tax burden.

Thus

\[
\boxed{
\text{tax incidence}
}
\]

is an elasticity problem.

Legal assignment of a tax does not by itself determine economic incidence.

%%%%%%%%%%%%%%%%%%%%%%%%%%%%%%%%%%%%%%%%%%%%%%%%%%%%%%%%%%%%
\subsection{Subsidies}
\label{subsec:economic-subsidies}
%%%%%%%%%%%%%%%%%%%%%%%%%%%%%%%%%%%%%%%%%%%%%%%%%%%%%%%%%%%%

A per-unit subsidy creates the opposite wedge.

It may increase equilibrium quantity while transferring resources to buyers,
sellers, or both depending on elasticities.

Like taxation, its distributional effect depends on market responsiveness.

%%%%%%%%%%%%%%%%%%%%%%%%%%%%%%%%%%%%%%%%%%%%%%%%%%%%%%%%%%%%
\subsection{Price Ceilings}
\label{subsec:price-ceilings}
%%%%%%%%%%%%%%%%%%%%%%%%%%%%%%%%%%%%%%%%%%%%%%%%%%%%%%%%%%%%

A binding price ceiling satisfies

\[
\boxed{
P_{\max}<P^*.
}
\]

At that price,

\[
Q_D(P_{\max})
>
Q_S(P_{\max}),
\]

creating excess demand in the simple model.

%%%%%%%%%%%%%%%%%%%%%%%%%%%%%%%%%%%%%%%%%%%%%%%%%%%%%%%%%%%%
\subsection{Price Floors}
\label{subsec:price-floors}
%%%%%%%%%%%%%%%%%%%%%%%%%%%%%%%%%%%%%%%%%%%%%%%%%%%%%%%%%%%%

A binding price floor satisfies

\[
\boxed{
P_{\min}>P^*.
}
\]

Then

\[
Q_S(P_{\min})
>
Q_D(P_{\min}),
\]

creating excess supply in the basic model.

%%%%%%%%%%%%%%%%%%%%%%%%%%%%%%%%%%%%%%%%%%%%%%%%%%%%%%%%%%%%
\subsection{Monopoly}
\label{subsec:monopoly}
%%%%%%%%%%%%%%%%%%%%%%%%%%%%%%%%%%%%%%%%%%%%%%%%%%%%%%%%%%%%

A monopolist recognizes that price depends on quantity.

Suppose inverse demand is

\[
P=P(q).
\]

Revenue is

\[
\boxed{
R(q)
=
P(q)q.
}
\]

Profit is

\[
\boxed{
\pi(q)
=
P(q)q-C(q).
}
\]

%%%%%%%%%%%%%%%%%%%%%%%%%%%%%%%%%%%%%%%%%%%%%%%%%%%%%%%%%%%%
\subsection{Marginal Revenue}
\label{subsec:monopoly-marginal-revenue}
%%%%%%%%%%%%%%%%%%%%%%%%%%%%%%%%%%%%%%%%%%%%%%%%%%%%%%%%%%%%

Differentiate revenue:

\[
R'(q)
=
P(q)
+
qP'(q).
\]

Therefore,

\[
\boxed{
MR(q)
=
P(q)
+
qP'(q).
}
\]

For downward-sloping demand,

\[
P'(q)<0,
\]

so marginal revenue is below price.

%%%%%%%%%%%%%%%%%%%%%%%%%%%%%%%%%%%%%%%%%%%%%%%%%%%%%%%%%%%%
\subsection{Monopoly Optimality Condition}
\label{subsec:monopoly-optimality}
%%%%%%%%%%%%%%%%%%%%%%%%%%%%%%%%%%%%%%%%%%%%%%%%%%%%%%%%%%%%

An interior monopoly optimum satisfies

\[
\boxed{
MR(q^*)=MC(q^*).
}
\]

The monopolist then uses the demand curve to determine price

\[
\boxed{
P^*=P(q^*).
}
\]

%%%%%%%%%%%%%%%%%%%%%%%%%%%%%%%%%%%%%%%%%%%%%%%%%%%%%%%%%%%%
\subsection{Worked Example: Monopoly}
\label{subsec:worked-monopoly}
%%%%%%%%%%%%%%%%%%%%%%%%%%%%%%%%%%%%%%%%%%%%%%%%%%%%%%%%%%%%

\begin{workedexamplebox}{Linear Inverse Demand}

Suppose

\[
P(q)=100-q
\]

and

\[
C(q)=20q.
\]

Revenue is

\[
R(q)
=
100q-q^2.
\]

Thus

\[
MR
=
100-2q.
\]

Marginal cost is

\[
MC=20.
\]

Set

\[
100-2q=20.
\]

Then

\[
q^*=40.
\]

The monopoly price is

\[
P^*
=
100-40
=
60.
\]

\begin{answerbox}
\[
\boxed{
q^*=40,
\qquad
P^*=60.
}
\]
\end{answerbox}

\end{workedexamplebox}

%%%%%%%%%%%%%%%%%%%%%%%%%%%%%%%%%%%%%%%%%%%%%%%%%%%%%%%%%%%%
\subsection{Price Discrimination}
\label{subsec:price-discrimination}
%%%%%%%%%%%%%%%%%%%%%%%%%%%%%%%%%%%%%%%%%%%%%%%%%%%%%%%%%%%%

Price discrimination occurs when a seller charges different effective prices
to different buyers or units for reasons beyond simple cost differences.

Mathematical models examine how segmentation, information, and demand
elasticity affect optimal prices.

%%%%%%%%%%%%%%%%%%%%%%%%%%%%%%%%%%%%%%%%%%%%%%%%%%%%%%%%%%%%
\subsection{Oligopoly}
\label{subsec:oligopoly}
%%%%%%%%%%%%%%%%%%%%%%%%%%%%%%%%%%%%%%%%%%%%%%%%%%%%%%%%%%%%

An oligopoly contains a small number of strategically interacting firms.

Each firm's optimal action depends on expectations about competitors.

Game theory is therefore central to oligopoly modeling.

%%%%%%%%%%%%%%%%%%%%%%%%%%%%%%%%%%%%%%%%%%%%%%%%%%%%%%%%%%%%
\subsection{Cournot Competition}
\label{subsec:cournot-competition}
%%%%%%%%%%%%%%%%%%%%%%%%%%%%%%%%%%%%%%%%%%%%%%%%%%%%%%%%%%%%

In Cournot competition, firms choose quantities.

For two firms,

\[
\boxed{
Q=q_1+q_2.
}
\]

Firm \(i\) solves

\[
\boxed{
\max_{q_i}
\quad
P(q_1+q_2)q_i-C_i(q_i).
}
\]

Its optimal quantity depends on the competitor's quantity.

%%%%%%%%%%%%%%%%%%%%%%%%%%%%%%%%%%%%%%%%%%%%%%%%%%%%%%%%%%%%
\subsection{Reaction Functions}
\label{subsec:reaction-functions}
%%%%%%%%%%%%%%%%%%%%%%%%%%%%%%%%%%%%%%%%%%%%%%%%%%%%%%%%%%%%

A reaction function gives one player's best response to another player's
choice.

For firm \(1\),

\[
\boxed{
q_1
=
R_1(q_2).
}
\]

For firm \(2\),

\[
\boxed{
q_2
=
R_2(q_1).
}
\]

A Cournot equilibrium occurs at an intersection.

%%%%%%%%%%%%%%%%%%%%%%%%%%%%%%%%%%%%%%%%%%%%%%%%%%%%%%%%%%%%
\subsection{Worked Example: Cournot Duopoly}
\label{subsec:worked-cournot}
%%%%%%%%%%%%%%%%%%%%%%%%%%%%%%%%%%%%%%%%%%%%%%%%%%%%%%%%%%%%

\begin{workedexamplebox}{Symmetric Duopoly}

Suppose

\[
P(Q)=a-bQ,
\]

with

\[
Q=q_1+q_2
\]

and zero marginal cost.

Firm \(1\)'s profit is

\[
\pi_1
=
(a-bq_1-bq_2)q_1.
\]

Differentiate:

\[
\frac{\partial\pi_1}{\partial q_1}
=
a-2bq_1-bq_2.
\]

Thus

\[
q_1
=
\frac{a-bq_2}{2b}.
\]

By symmetry,

\[
q_2
=
\frac{a-bq_1}{2b}.
\]

Solving gives

\[
\begin{answerbox}
\[
\boxed{
q_1^*=q_2^*=\frac{a}{3b}.
}
\]
\end{answerbox}

\end{workedexamplebox}

%%%%%%%%%%%%%%%%%%%%%%%%%%%%%%%%%%%%%%%%%%%%%%%%%%%%%%%%%%%%
\subsection{Bertrand Competition}
\label{subsec:bertrand-competition}
%%%%%%%%%%%%%%%%%%%%%%%%%%%%%%%%%%%%%%%%%%%%%%%%%%%%%%%%%%%%

In Bertrand competition, firms choose prices rather than quantities.

The outcome depends strongly on assumptions about:

\[
\boxed{
\text{product differentiation},
}
\]

\[
\boxed{
\text{capacity},
}
\]

and

\[
\boxed{
\text{cost structure}.
}
\]

For identical products and constant marginal costs under classical
assumptions, price competition can drive price toward marginal cost.

%%%%%%%%%%%%%%%%%%%%%%%%%%%%%%%%%%%%%%%%%%%%%%%%%%%%%%%%%%%%
\subsection{Stackelberg Competition}
\label{subsec:stackelberg-competition}
%%%%%%%%%%%%%%%%%%%%%%%%%%%%%%%%%%%%%%%%%%%%%%%%%%%%%%%%%%%%

In Stackelberg competition, one firm acts first and another responds.

The leader anticipates the follower's reaction function.

The model is solved by backward induction.

Thus sequential timing changes strategic outcomes.

%%%%%%%%%%%%%%%%%%%%%%%%%%%%%%%%%%%%%%%%%%%%%%%%%%%%%%%%%%%%
\subsection{Nash Equilibrium}
\label{subsec:nash-equilibrium-economics}
%%%%%%%%%%%%%%%%%%%%%%%%%%%%%%%%%%%%%%%%%%%%%%%%%%%%%%%%%%%%

A strategy profile

\[
\boxed{
(s_1^*,\ldots,s_n^*)
}
\]

is a Nash equilibrium if no player can improve by unilateral deviation:

\[
\boxed{
u_i(s_i^*,s_{-i}^*)
\geq
u_i(s_i,s_{-i}^*)
}
\]

for every feasible \(s_i\).

%%%%%%%%%%%%%%%%%%%%%%%%%%%%%%%%%%%%%%%%%%%%%%%%%%%%%%%%%%%%
\subsection{Dominant Strategies}
\label{subsec:dominant-strategies}
%%%%%%%%%%%%%%%%%%%%%%%%%%%%%%%%%%%%%%%%%%%%%%%%%%%%%%%%%%%%

A strategy is dominant if it gives at least as high a payoff as every other
strategy regardless of opponents' actions.

A dominant-strategy equilibrium is therefore especially strong.

Not every game has dominant strategies.

%%%%%%%%%%%%%%%%%%%%%%%%%%%%%%%%%%%%%%%%%%%%%%%%%%%%%%%%%%%%
\subsection{Mixed Strategies}
\label{subsec:mixed-strategies-economics}
%%%%%%%%%%%%%%%%%%%%%%%%%%%%%%%%%%%%%%%%%%%%%%%%%%%%%%%%%%%%

A mixed strategy assigns probabilities to pure strategies.

For finite games, mixed strategies enlarge the strategy space and guarantee
existence of Nash equilibrium under standard assumptions.

Expected payoffs are computed using probability.

%%%%%%%%%%%%%%%%%%%%%%%%%%%%%%%%%%%%%%%%%%%%%%%%%%%%%%%%%%%%
\subsection{Repeated Games}
\label{subsec:repeated-games-economics}
%%%%%%%%%%%%%%%%%%%%%%%%%%%%%%%%%%%%%%%%%%%%%%%%%%%%%%%%%%%%

When a game is repeated, players may condition current actions on past
behavior.

This can support outcomes that differ from the one-shot equilibrium.

Discount factors determine how strongly players value future consequences.

%%%%%%%%%%%%%%%%%%%%%%%%%%%%%%%%%%%%%%%%%%%%%%%%%%%%%%%%%%%%
\subsection{Public Goods}
\label{subsec:public-goods}
%%%%%%%%%%%%%%%%%%%%%%%%%%%%%%%%%%%%%%%%%%%%%%%%%%%%%%%%%%%%

A public good is commonly characterized as:

\[
\boxed{
\text{nonrival}
}
\]

and

\[
\boxed{
\text{nonexcludable}.
}
\]

Individual incentives to contribute may differ from socially efficient
provision.

This creates a strategic collective-action problem.

%%%%%%%%%%%%%%%%%%%%%%%%%%%%%%%%%%%%%%%%%%%%%%%%%%%%%%%%%%%%
\subsection{Externalities}
\label{subsec:economic-externalities}
%%%%%%%%%%%%%%%%%%%%%%%%%%%%%%%%%%%%%%%%%%%%%%%%%%%%%%%%%%%%

An externality occurs when one agent's action affects another agent's welfare
outside the modeled market price mechanism.

Examples include:

\[
\boxed{
\text{pollution}
}
\]

and

\[
\boxed{
\text{knowledge spillovers}.
}
\]

The private optimum may then differ from the social optimum.

%%%%%%%%%%%%%%%%%%%%%%%%%%%%%%%%%%%%%%%%%%%%%%%%%%%%%%%%%%%%
\subsection{Pigouvian Taxes}
\label{subsec:pigouvian-taxes}
%%%%%%%%%%%%%%%%%%%%%%%%%%%%%%%%%%%%%%%%%%%%%%%%%%%%%%%%%%%%

A Pigouvian tax attempts to make private decision makers internalize a
negative externality.

In a simple model, the tax is chosen to reflect marginal external damage at
the socially efficient quantity.

Thus

\[
\boxed{
\text{private marginal cost}
+
\text{external marginal cost}
}
\]

becomes relevant to the social decision.

%%%%%%%%%%%%%%%%%%%%%%%%%%%%%%%%%%%%%%%%%%%%%%%%%%%%%%%%%%%%
\subsection{Common-Pool Resources}
\label{subsec:common-pool-resources}
%%%%%%%%%%%%%%%%%%%%%%%%%%%%%%%%%%%%%%%%%%%%%%%%%%%%%%%%%%%%

A common-pool resource is difficult to exclude users from but is rival in
consumption.

Examples include:

\[
\boxed{
\text{fisheries},
}
\]

\[
\boxed{
\text{groundwater},
}
\]

and

\[
\boxed{
\text{shared grazing land}.
}
\]

Individual incentives may produce overuse.

%%%%%%%%%%%%%%%%%%%%%%%%%%%%%%%%%%%%%%%%%%%%%%%%%%%%%%%%%%%%
\subsection{Expected Utility}
\label{subsec:expected-utility-economics}
%%%%%%%%%%%%%%%%%%%%%%%%%%%%%%%%%%%%%%%%%%%%%%%%%%%%%%%%%%%%

Suppose wealth outcome \(W_i\) occurs with probability \(p_i\).

Expected utility is

\[
\boxed{
\mathbb{E}[u(W)]
=
\sum_i
p_i u(W_i).
}
\]

A decision maker chooses among lotteries by comparing expected utility under
the model.

%%%%%%%%%%%%%%%%%%%%%%%%%%%%%%%%%%%%%%%%%%%%%%%%%%%%%%%%%%%%
\subsection{Risk Aversion}
\label{subsec:risk-aversion}
%%%%%%%%%%%%%%%%%%%%%%%%%%%%%%%%%%%%%%%%%%%%%%%%%%%%%%%%%%%%

A twice-differentiable utility function represents risk aversion when

\[
\boxed{
u''(W)<0.
}
\]

Thus utility is concave in wealth.

By Jensen's inequality,

\[
\boxed{
\mathbb{E}[u(W)]
\leq
u(\mathbb{E}[W]).
}
\]

%%%%%%%%%%%%%%%%%%%%%%%%%%%%%%%%%%%%%%%%%%%%%%%%%%%%%%%%%%%%
\subsection{Certainty Equivalent}
\label{subsec:certainty-equivalent}
%%%%%%%%%%%%%%%%%%%%%%%%%%%%%%%%%%%%%%%%%%%%%%%%%%%%%%%%%%%%

The certainty equivalent \(CE\) satisfies

\[
\boxed{
u(CE)
=
\mathbb{E}[u(W)].
}
\]

For a risk-averse agent,

\[
\boxed{
CE
\leq
\mathbb{E}[W].
}
\]

%%%%%%%%%%%%%%%%%%%%%%%%%%%%%%%%%%%%%%%%%%%%%%%%%%%%%%%%%%%%
\subsection{Risk Premium}
\label{subsec:risk-premium-economics}
%%%%%%%%%%%%%%%%%%%%%%%%%%%%%%%%%%%%%%%%%%%%%%%%%%%%%%%%%%%%

The risk premium is

\[
\boxed{
\pi
=
\mathbb{E}[W]-CE.
}
\]

For a risk-averse decision maker,

\[
\boxed{
\pi\geq0.
}
\]

It measures how much expected wealth the agent would sacrifice to eliminate
risk.

%%%%%%%%%%%%%%%%%%%%%%%%%%%%%%%%%%%%%%%%%%%%%%%%%%%%%%%%%%%%
\subsection{Arrow--Pratt Risk Aversion}
\label{subsec:arrow-pratt}
%%%%%%%%%%%%%%%%%%%%%%%%%%%%%%%%%%%%%%%%%%%%%%%%%%%%%%%%%%%%

Absolute risk aversion is

\[
\boxed{
A(W)
=
-
\frac{u''(W)}{u'(W)}.
}
\]

Relative risk aversion is

\[
\boxed{
R(W)
=
-
W
\frac{u''(W)}{u'(W)}.
}
\]

These measures quantify local curvature of utility.

%%%%%%%%%%%%%%%%%%%%%%%%%%%%%%%%%%%%%%%%%%%%%%%%%%%%%%%%%%%%
\subsection{Intertemporal Choice}
\label{subsec:intertemporal-choice}
%%%%%%%%%%%%%%%%%%%%%%%%%%%%%%%%%%%%%%%%%%%%%%%%%%%%%%%%%%%%

Economic decisions often trade present and future consumption.

A simple objective is

\[
\boxed{
U
=
u(c_0)
+
\beta u(c_1),
}
\]

where

\[
0<\beta\leq1
\]

is a discount factor.

Smaller \(\beta\) places less weight on future utility.

%%%%%%%%%%%%%%%%%%%%%%%%%%%%%%%%%%%%%%%%%%%%%%%%%%%%%%%%%%%%
\subsection{Consumption-Saving Constraint}
\label{subsec:consumption-saving}
%%%%%%%%%%%%%%%%%%%%%%%%%%%%%%%%%%%%%%%%%%%%%%%%%%%%%%%%%%%%

Suppose current resources are \(Y_0\), future resources are \(Y_1\), and the
interest rate is \(r\).

An intertemporal budget can be written

\[
\boxed{
c_0
+
\frac{c_1}{1+r}
=
Y_0
+
\frac{Y_1}{1+r}.
}
\]

The consumer chooses consumption across time.

%%%%%%%%%%%%%%%%%%%%%%%%%%%%%%%%%%%%%%%%%%%%%%%%%%%%%%%%%%%%
\subsection{Euler Equation for Consumption}
\label{subsec:consumption-euler-equation}
%%%%%%%%%%%%%%%%%%%%%%%%%%%%%%%%%%%%%%%%%%%%%%%%%%%%%%%%%%%%

For

\[
\max
\quad
u(c_0)+\beta u(c_1)
\]

subject to the intertemporal budget, an interior optimum satisfies

\[
\boxed{
u'(c_0)
=
\beta(1+r)u'(c_1).
}
\]

This balances marginal utility across time.

%%%%%%%%%%%%%%%%%%%%%%%%%%%%%%%%%%%%%%%%%%%%%%%%%%%%%%%%%%%%
\subsection{Dynamic Optimization}
\label{subsec:dynamic-economic-optimization}
%%%%%%%%%%%%%%%%%%%%%%%%%%%%%%%%%%%%%%%%%%%%%%%%%%%%%%%%%%%%

Many economic problems have state variables evolving through time.

A general discrete-time system is

\[
\boxed{
x_{t+1}
=
f(x_t,u_t).
}
\]

The decision maker chooses controls \(u_t\) to maximize

\[
\boxed{
\sum_{t=0}^{\infty}
\beta^t
r(x_t,u_t).
}
\]

%%%%%%%%%%%%%%%%%%%%%%%%%%%%%%%%%%%%%%%%%%%%%%%%%%%%%%%%%%%%
\subsection{Bellman Equation}
\label{subsec:economic-bellman-equation}
%%%%%%%%%%%%%%%%%%%%%%%%%%%%%%%%%%%%%%%%%%%%%%%%%%%%%%%%%%%%

The value function satisfies

\[
\boxed{
V(x)
=
\max_u
\left\{
r(x,u)
+
\beta V(f(x,u))
\right\}.
}
\]

This is the Bellman equation.

It converts an infinite sequence problem into recursive optimization.

%%%%%%%%%%%%%%%%%%%%%%%%%%%%%%%%%%%%%%%%%%%%%%%%%%%%%%%%%%%%
\subsection{Economic Steady States}
\label{subsec:economic-steady-states}
%%%%%%%%%%%%%%%%%%%%%%%%%%%%%%%%%%%%%%%%%%%%%%%%%%%%%%%%%%%%

For a dynamic system

\[
x_{t+1}
=
F(x_t),
\]

a steady state satisfies

\[
\boxed{
x^*
=
F(x^*).
}
\]

In continuous time,

\[
\dot{x}
=
f(x),
\]

a steady state satisfies

\[
\boxed{
f(x^*)=0.
}
\]

%%%%%%%%%%%%%%%%%%%%%%%%%%%%%%%%%%%%%%%%%%%%%%%%%%%%%%%%%%%%
\subsection{Local Stability}
\label{subsec:economic-local-stability}
%%%%%%%%%%%%%%%%%%%%%%%%%%%%%%%%%%%%%%%%%%%%%%%%%%%%%%%%%%%%

For a one-dimensional discrete model,

\[
x_{t+1}
=
F(x_t),
\]

a fixed point \(x^*\) is locally asymptotically stable when

\[
\boxed{
|F'(x^*)|<1.
}
\]

For continuous-time

\[
\dot{x}=f(x),
\]

local asymptotic stability in one dimension requires

\[
\boxed{
f'(x^*)<0.
}
\]

%%%%%%%%%%%%%%%%%%%%%%%%%%%%%%%%%%%%%%%%%%%%%%%%%%%%%%%%%%%%
\subsection{Solow Growth Model}
\label{subsec:solow-growth-model}
%%%%%%%%%%%%%%%%%%%%%%%%%%%%%%%%%%%%%%%%%%%%%%%%%%%%%%%%%%%%

A basic Solow model for capital per worker is

\[
\boxed{
\dot{k}
=
sf(k)
-
(n+\delta)k.
}
\]

Here:

\[
s
=
\text{saving rate},
\]

\[
n
=
\text{population growth rate},
\]

and

\[
\delta
=
\text{depreciation rate}.
\]

%%%%%%%%%%%%%%%%%%%%%%%%%%%%%%%%%%%%%%%%%%%%%%%%%%%%%%%%%%%%
\subsection{Solow Steady State}
\label{subsec:solow-steady-state}
%%%%%%%%%%%%%%%%%%%%%%%%%%%%%%%%%%%%%%%%%%%%%%%%%%%%%%%%%%%%

A positive steady state satisfies

\[
\boxed{
sf(k^*)
=
(n+\delta)k^*.
}
\]

Investment per worker equals effective capital dilution.

The dynamics determine whether capital converges toward this value.

%%%%%%%%%%%%%%%%%%%%%%%%%%%%%%%%%%%%%%%%%%%%%%%%%%%%%%%%%%%%
\subsection{Worked Example: Solow Steady State}
\label{subsec:worked-solow}
%%%%%%%%%%%%%%%%%%%%%%%%%%%%%%%%%%%%%%%%%%%%%%%%%%%%%%%%%%%%

\begin{workedexamplebox}{Square-Root Production}

Suppose

\[
f(k)=\sqrt{k}.
\]

Then the steady-state equation is

\[
s\sqrt{k}
=
(n+\delta)k.
\]

For \(k>0\),

\[
s
=
(n+\delta)\sqrt{k}.
\]

Thus

\[
\sqrt{k^*}
=
\frac{s}{n+\delta}.
\]

Therefore,

\begin{answerbox}
\[
\boxed{
k^*
=
\left(
\frac{s}{n+\delta}
\right)^2.
}
\]
\end{answerbox}

\end{workedexamplebox}

%%%%%%%%%%%%%%%%%%%%%%%%%%%%%%%%%%%%%%%%%%%%%%%%%%%%%%%%%%%%
\subsection{Optimal Growth}
\label{subsec:optimal-growth}
%%%%%%%%%%%%%%%%%%%%%%%%%%%%%%%%%%%%%%%%%%%%%%%%%%%%%%%%%%%%

Optimal growth models choose consumption and investment over time to maximize
intertemporal welfare.

A stylized problem is

\[
\boxed{
\max_{c(t)}
\int_0^\infty
e^{-\rho t}
u(c(t))
\,dt
}
\]

subject to capital dynamics

\[
\boxed{
\dot{k}
=
f(k)-c-\delta k.
}
\]

This is a continuous-time optimal control problem.

%%%%%%%%%%%%%%%%%%%%%%%%%%%%%%%%%%%%%%%%%%%%%%%%%%%%%%%%%%%%
\subsection{Input--Output Economics}
\label{subsec:input-output-economics}
%%%%%%%%%%%%%%%%%%%%%%%%%%%%%%%%%%%%%%%%%%%%%%%%%%%%%%%%%%%%

Input--output models describe how industries use outputs of other industries
as production inputs.

Let

\[
\boxed{
A
}
\]

be the matrix of technical coefficients.

If \(\mathbf{x}\) is total production and \(\mathbf{d}\) final demand, then

\[
\boxed{
\mathbf{x}
=
A\mathbf{x}
+
\mathbf{d}.
}
\]

%%%%%%%%%%%%%%%%%%%%%%%%%%%%%%%%%%%%%%%%%%%%%%%%%%%%%%%%%%%%
\subsection{Leontief Model}
\label{subsec:leontief-model}
%%%%%%%%%%%%%%%%%%%%%%%%%%%%%%%%%%%%%%%%%%%%%%%%%%%%%%%%%%%%

Rearranging,

\[
\boxed{
(I-A)\mathbf{x}
=
\mathbf{d}.
}
\]

If \(I-A\) is invertible,

\[
\boxed{
\mathbf{x}
=
(I-A)^{-1}\mathbf{d}.
}
\]

The matrix

\[
\boxed{
(I-A)^{-1}
}
\]

is the Leontief inverse.

%%%%%%%%%%%%%%%%%%%%%%%%%%%%%%%%%%%%%%%%%%%%%%%%%%%%%%%%%%%%
\subsection{Worked Example: Two-Sector Leontief Model}
\label{subsec:worked-leontief}
%%%%%%%%%%%%%%%%%%%%%%%%%%%%%%%%%%%%%%%%%%%%%%%%%%%%%%%%%%%%

\begin{workedexamplebox}{Production Requirements}

Suppose

\[
A
=
\begin{bmatrix}
0.2 & 0.1\\
0.3 & 0.2
\end{bmatrix},
\qquad
\mathbf{d}
=
\begin{bmatrix}
100\\
80
\end{bmatrix}.
\]

Then

\[
I-A
=
\begin{bmatrix}
0.8 & -0.1\\
-0.3 & 0.8
\end{bmatrix}.
\]

The production vector satisfies

\[
\mathbf{x}
=
(I-A)^{-1}\mathbf{d}.
\]

Since

\[
\det(I-A)
=
0.64-0.03
=
0.61,
\]

\[
(I-A)^{-1}
=
\frac{1}{0.61}
\begin{bmatrix}
0.8 & 0.1\\
0.3 & 0.8
\end{bmatrix}.
\]

Therefore,

\[
\mathbf{x}
=
\frac{1}{0.61}
\begin{bmatrix}
88\\
94
\end{bmatrix}.
\]

Hence

\begin{answerbox}
\[
\boxed{
\mathbf{x}
\approx
\begin{bmatrix}
144.26\\
154.10
\end{bmatrix}.
}
\]
\end{answerbox}

\end{workedexamplebox}

%%%%%%%%%%%%%%%%%%%%%%%%%%%%%%%%%%%%%%%%%%%%%%%%%%%%%%%%%%%%
\subsection{General Equilibrium}
\label{subsec:general-equilibrium}
%%%%%%%%%%%%%%%%%%%%%%%%%%%%%%%%%%%%%%%%%%%%%%%%%%%%%%%%%%%%

General equilibrium studies simultaneous equilibrium across many markets.

Prices affect:

\[
\boxed{
\text{consumer demand},
}
\]

\[
\boxed{
\text{firm supply},
}
\]

and

\[
\boxed{
\text{resource allocation}.
}
\]

An equilibrium price vector makes aggregate excess demand compatible with
market clearing.

%%%%%%%%%%%%%%%%%%%%%%%%%%%%%%%%%%%%%%%%%%%%%%%%%%%%%%%%%%%%
\subsection{Excess Demand}
\label{subsec:excess-demand}
%%%%%%%%%%%%%%%%%%%%%%%%%%%%%%%%%%%%%%%%%%%%%%%%%%%%%%%%%%%%

Let

\[
\mathbf{z}(\mathbf{p})
\]

be aggregate excess demand.

A market-clearing equilibrium satisfies

\[
\boxed{
\mathbf{z}(\mathbf{p}^*)=\mathbf{0},
}
\]

subject to an appropriate normalization of the price vector.

%%%%%%%%%%%%%%%%%%%%%%%%%%%%%%%%%%%%%%%%%%%%%%%%%%%%%%%%%%%%
\subsection{Walras' Law}
\label{subsec:walras-law}
%%%%%%%%%%%%%%%%%%%%%%%%%%%%%%%%%%%%%%%%%%%%%%%%%%%%%%%%%%%%

In standard general-equilibrium formulations,

\[
\boxed{
\mathbf{p}^T
\mathbf{z}(\mathbf{p})
=
0.
}
\]

This is Walras' law.

It means the value of aggregate excess demand sums to zero under the model's
budget constraints.

%%%%%%%%%%%%%%%%%%%%%%%%%%%%%%%%%%%%%%%%%%%%%%%%%%%%%%%%%%%%
\subsection{Fixed Points and Equilibrium}
\label{subsec:fixed-points-economic-equilibrium}
%%%%%%%%%%%%%%%%%%%%%%%%%%%%%%%%%%%%%%%%%%%%%%%%%%%%%%%%%%%%

Existence of economic equilibrium is often formulated using fixed-point
theorems.

A fixed point satisfies

\[
\boxed{
x^*
=
F(x^*).
}
\]

Thus equilibrium theory connects economics with topology and functional
analysis.

%%%%%%%%%%%%%%%%%%%%%%%%%%%%%%%%%%%%%%%%%%%%%%%%%%%%%%%%%%%%
\subsection{Auctions}
\label{subsec:auction-models}
%%%%%%%%%%%%%%%%%%%%%%%%%%%%%%%%%%%%%%%%%%%%%%%%%%%%%%%%%%%%

An auction allocates an object or resource using bids.

Mathematical auction models specify:

\[
\boxed{
\text{valuation distributions},
}
\]

\[
\boxed{
\text{bidding strategies},
}
\]

and

\[
\boxed{
\text{allocation and payment rules}.
}
\]

Game theory and probability are central to auction analysis.

%%%%%%%%%%%%%%%%%%%%%%%%%%%%%%%%%%%%%%%%%%%%%%%%%%%%%%%%%%%%
\subsection{First-Price Auctions}
\label{subsec:first-price-auctions}
%%%%%%%%%%%%%%%%%%%%%%%%%%%%%%%%%%%%%%%%%%%%%%%%%%%%%%%%%%%%

In a first-price sealed-bid auction, the highest bidder wins and pays their
own bid.

A bidder balances:

\[
\boxed{
\text{higher probability of winning}
}
\]

against

\[
\boxed{
\text{lower surplus if the bid is high}.
}
\]

Optimal bidding therefore becomes a strategic optimization problem.

%%%%%%%%%%%%%%%%%%%%%%%%%%%%%%%%%%%%%%%%%%%%%%%%%%%%%%%%%%%%
\subsection{Second-Price Auctions}
\label{subsec:second-price-auctions}
%%%%%%%%%%%%%%%%%%%%%%%%%%%%%%%%%%%%%%%%%%%%%%%%%%%%%%%%%%%%

In a second-price sealed-bid auction, the highest bidder wins but pays the
second-highest bid.

Under standard private-value assumptions, truthful bidding is a weakly
dominant strategy.

This is an important example of mechanism design.

%%%%%%%%%%%%%%%%%%%%%%%%%%%%%%%%%%%%%%%%%%%%%%%%%%%%%%%%%%%%
\subsection{Mechanism Design}
\label{subsec:mechanism-design}
%%%%%%%%%%%%%%%%%%%%%%%%%%%%%%%%%%%%%%%%%%%%%%%%%%%%%%%%%%%%

Mechanism design asks how rules should be constructed so that strategic
behavior produces desired outcomes.

The designer chooses:

\[
\boxed{
\text{allocation rules}
}
\]

and

\[
\boxed{
\text{payment rules}
}
\]

subject to incentive and participation constraints.

%%%%%%%%%%%%%%%%%%%%%%%%%%%%%%%%%%%%%%%%%%%%%%%%%%%%%%%%%%%%
\subsection{Incentive Compatibility}
\label{subsec:incentive-compatibility}
%%%%%%%%%%%%%%%%%%%%%%%%%%%%%%%%%%%%%%%%%%%%%%%%%%%%%%%%%%%%

A mechanism is incentive compatible when agents prefer to follow the intended
reporting or strategy rule.

Schematically,

\[
\boxed{
u_i(\text{truthful})
\geq
u_i(\text{deviation})
}
\]

for the relevant information and strategic environment.

%%%%%%%%%%%%%%%%%%%%%%%%%%%%%%%%%%%%%%%%%%%%%%%%%%%%%%%%%%%%
\subsection{Social Choice}
\label{subsec:social-choice}
%%%%%%%%%%%%%%%%%%%%%%%%%%%%%%%%%%%%%%%%%%%%%%%%%%%%%%%%%%%%

Social choice studies how individual preferences are aggregated into
collective decisions.

A social choice rule maps preference profiles to:

\[
\boxed{
\text{rankings}
}
\]

or

\[
\boxed{
\text{outcomes}.
}
\]

Different fairness and consistency conditions may conflict.

%%%%%%%%%%%%%%%%%%%%%%%%%%%%%%%%%%%%%%%%%%%%%%%%%%%%%%%%%%%%
\subsection{Econometrics Connection}
\label{subsec:econometrics-connection}
%%%%%%%%%%%%%%%%%%%%%%%%%%%%%%%%%%%%%%%%%%%%%%%%%%%%%%%%%%%%

Economic theory proposes relationships among variables.

Econometrics uses data to estimate and test those relationships.

A simple linear model is

\[
\boxed{
Y
=
X\beta+\varepsilon.
}
\]

The mathematical structure is statistical, but interpretation depends on the
economic model.

%%%%%%%%%%%%%%%%%%%%%%%%%%%%%%%%%%%%%%%%%%%%%%%%%%%%%%%%%%%%
\subsection{Structural Models}
\label{subsec:structural-economic-models}
%%%%%%%%%%%%%%%%%%%%%%%%%%%%%%%%%%%%%%%%%%%%%%%%%%%%%%%%%%%%

A structural model attempts to represent behavioral or technological
relationships derived from economic assumptions.

Its parameters may have direct interpretations such as:

\[
\boxed{
\text{preferences},
}
\]

\[
\boxed{
\text{technology},
}
\]

or

\[
\boxed{
\text{policy response}.
}
\]

%%%%%%%%%%%%%%%%%%%%%%%%%%%%%%%%%%%%%%%%%%%%%%%%%%%%%%%%%%%%
\subsection{Reduced-Form Models}
\label{subsec:reduced-form-economic-models}
%%%%%%%%%%%%%%%%%%%%%%%%%%%%%%%%%%%%%%%%%%%%%%%%%%%%%%%%%%%%

A reduced-form model expresses endogenous outcomes directly as functions of
observable exogenous variables and disturbances.

Reduced forms may be easier to estimate but can contain less structural
interpretation.

%%%%%%%%%%%%%%%%%%%%%%%%%%%%%%%%%%%%%%%%%%%%%%%%%%%%%%%%%%%%
\subsection{Identification}
\label{subsec:economic-identification}
%%%%%%%%%%%%%%%%%%%%%%%%%%%%%%%%%%%%%%%%%%%%%%%%%%%%%%%%%%%%

Identification asks whether model parameters can be uniquely determined from
the probability distribution of observable data.

If two different parameter values imply the same observable distribution,
the parameter is not identified.

Thus

\[
\boxed{
\text{more data}
}
\]

does not solve a fundamentally unidentified model.

%%%%%%%%%%%%%%%%%%%%%%%%%%%%%%%%%%%%%%%%%%%%%%%%%%%%%%%%%%%%
\subsection{Simultaneity}
\label{subsec:economic-simultaneity}
%%%%%%%%%%%%%%%%%%%%%%%%%%%%%%%%%%%%%%%%%%%%%%%%%%%%%%%%%%%%

Supply and demand illustrate simultaneity.

Observed equilibrium price and quantity are jointly determined.

Regressing quantity on price without addressing this joint determination can
fail to recover a demand relationship.

This is a structural identification problem.

%%%%%%%%%%%%%%%%%%%%%%%%%%%%%%%%%%%%%%%%%%%%%%%%%%%%%%%%%%%%
\subsection{Policy Optimization}
\label{subsec:economic-policy-optimization}
%%%%%%%%%%%%%%%%%%%%%%%%%%%%%%%%%%%%%%%%%%%%%%%%%%%%%%%%%%%%

A policymaker may solve

\[
\boxed{
\max_{\mathbf{u}}
\quad
W(\mathbf{x},\mathbf{u})
}
\]

subject to economic dynamics

\[
\boxed{
\dot{\mathbf{x}}
=
f(\mathbf{x},\mathbf{u})
}
\]

and policy constraints.

This connects economics with optimal control and operations research.

%%%%%%%%%%%%%%%%%%%%%%%%%%%%%%%%%%%%%%%%%%%%%%%%%%%%%%%%%%%%
\subsection{Economic Networks}
\label{subsec:economic-networks}
%%%%%%%%%%%%%%%%%%%%%%%%%%%%%%%%%%%%%%%%%%%%%%%%%%%%%%%%%%%%

Economic behavior often occurs through networks.

Examples include:

\[
\boxed{
\text{trade networks},
}
\]

\[
\boxed{
\text{production networks},
}
\]

\[
\boxed{
\text{labor networks},
}
\]

and

\[
\boxed{
\text{financial networks}.
}
\]

Network position can affect prices, productivity, risk, and diffusion.

%%%%%%%%%%%%%%%%%%%%%%%%%%%%%%%%%%%%%%%%%%%%%%%%%%%%%%%%%%%%
\subsection{Network Effects}
\label{subsec:network-effects-economics}
%%%%%%%%%%%%%%%%%%%%%%%%%%%%%%%%%%%%%%%%%%%%%%%%%%%%%%%%%%%%

A network effect occurs when the value or benefit of a good depends on how
many others use it or how they are connected.

A stylized utility might be

\[
\boxed{
u_i
=
v_i(x_i)
+
\alpha
\sum_jA_{ij}x_ix_j.
}
\]

Thus individual decisions depend on neighbors' behavior.

%%%%%%%%%%%%%%%%%%%%%%%%%%%%%%%%%%%%%%%%%%%%%%%%%%%%%%%%%%%%
\subsection{Multiple Equilibria}
\label{subsec:multiple-equilibria-economics}
%%%%%%%%%%%%%%%%%%%%%%%%%%%%%%%%%%%%%%%%%%%%%%%%%%%%%%%%%%%%

Nonlinear interactions or strategic complementarities can generate multiple
equilibria.

Thus the model may admit

\[
\boxed{
x_1^*,
\quad
x_2^*,
\quad
\ldots
}
\]

rather than a unique solution.

History, expectations, and coordination can then influence which equilibrium
is observed.

%%%%%%%%%%%%%%%%%%%%%%%%%%%%%%%%%%%%%%%%%%%%%%%%%%%%%%%%%%%%
\subsection{Expectations}
\label{subsec:economic-expectations}
%%%%%%%%%%%%%%%%%%%%%%%%%%%%%%%%%%%%%%%%%%%%%%%%%%%%%%%%%%%%

Economic decisions often depend on beliefs about the future.

For example,

\[
\boxed{
q_t
=
f
\left(
p_t,
\mathbb{E}_t[p_{t+1}]
\right).
}
\]

Expectations can therefore become endogenous components of dynamic models.

%%%%%%%%%%%%%%%%%%%%%%%%%%%%%%%%%%%%%%%%%%%%%%%%%%%%%%%%%%%%
\subsection{Rational Expectations}
\label{subsec:rational-expectations}
%%%%%%%%%%%%%%%%%%%%%%%%%%%%%%%%%%%%%%%%%%%%%%%%%%%%%%%%%%%%

In a rational-expectations model, agents' expectations are consistent with
the probability structure implied by the model.

This does not mean perfect foresight.

Random future outcomes may remain uncertain.

%%%%%%%%%%%%%%%%%%%%%%%%%%%%%%%%%%%%%%%%%%%%%%%%%%%%%%%%%%%%
\subsection{Model Calibration}
\label{subsec:economic-calibration}
%%%%%%%%%%%%%%%%%%%%%%%%%%%%%%%%%%%%%%%%%%%%%%%%%%%%%%%%%%%%

Calibration chooses parameter values so model outputs match selected
empirical targets.

If moments generated by a model are

\[
m_j(\theta),
\]

one may choose

\[
\boxed{
\theta
}
\]

to make them close to observed targets

\[
\hat m_j.
\]

%%%%%%%%%%%%%%%%%%%%%%%%%%%%%%%%%%%%%%%%%%%%%%%%%%%%%%%%%%%%
\subsection{Sensitivity Analysis}
\label{subsec:economic-sensitivity-analysis}
%%%%%%%%%%%%%%%%%%%%%%%%%%%%%%%%%%%%%%%%%%%%%%%%%%%%%%%%%%%%

Economic conclusions can depend strongly on parameters.

A sensitivity analysis varies:

\[
\boxed{
\text{elasticities},
}
\]

\[
\boxed{
\text{discount factors},
}
\]

\[
\boxed{
\text{growth rates},
}
\]

and

\[
\boxed{
\text{behavioral parameters}
}
\]

to determine whether conclusions remain robust.

%%%%%%%%%%%%%%%%%%%%%%%%%%%%%%%%%%%%%%%%%%%%%%%%%%%%%%%%%%%%
\subsection{Equilibrium Versus Prediction}
\label{subsec:equilibrium-versus-prediction}
%%%%%%%%%%%%%%%%%%%%%%%%%%%%%%%%%%%%%%%%%%%%%%%%%%%%%%%%%%%%

An equilibrium is a mathematical consistency condition.

It does not necessarily imply:

\[
\boxed{
\text{the economy reaches it quickly},
}
\]

\[
\boxed{
\text{the equilibrium is unique},
}
\]

or

\[
\boxed{
\text{the equilibrium is empirically accurate}.
}
\]

Dynamics and validation remain separate questions.

%%%%%%%%%%%%%%%%%%%%%%%%%%%%%%%%%%%%%%%%%%%%%%%%%%%%%%%%%%%%
\subsection{Positive and Normative Analysis}
\label{subsec:positive-normative}
%%%%%%%%%%%%%%%%%%%%%%%%%%%%%%%%%%%%%%%%%%%%%%%%%%%%%%%%%%%%

Positive economics asks what follows from specified behavioral and market
assumptions.

Normative economics evaluates outcomes using an objective such as welfare.

Mathematics can clarify both, but a normative objective requires value
judgments that are not determined by mathematics alone.

%%%%%%%%%%%%%%%%%%%%%%%%%%%%%%%%%%%%%%%%%%%%%%%%%%%%%%%%%%%%
\subsection{Mathematical Economics Workflow}
\label{subsec:mathematical-economics-workflow}
%%%%%%%%%%%%%%%%%%%%%%%%%%%%%%%%%%%%%%%%%%%%%%%%%%%%%%%%%%%%

A useful workflow is:

\begin{enumerate}
    \item define the economic question;
    \item identify agents and institutions;
    \item specify decision variables;
    \item distinguish endogenous and exogenous variables;
    \item define preferences or objectives;
    \item define technological constraints;
    \item define budget or resource constraints;
    \item specify market structure;
    \item identify strategic interactions;
    \item identify uncertainty;
    \item determine whether the model is static or dynamic;
    \item formulate optimization problems;
    \item derive first-order and second-order conditions;
    \item analyze constraints using KKT conditions when needed;
    \item derive demand, supply, or policy functions;
    \item impose equilibrium conditions;
    \item test existence and uniqueness when possible;
    \item analyze stability for dynamic equilibria;
    \item compute comparative statics;
    \item identify welfare implications;
    \item estimate or calibrate parameters;
    \item perform sensitivity analysis;
    \item compare model implications with data;
    \item distinguish descriptive conclusions from normative conclusions;
    \item communicate assumptions and limitations.
\end{enumerate}

%%%%%%%%%%%%%%%%%%%%%%%%%%%%%%%%%%%%%%%%%%%%%%%%%%%%%%%%%%%%
\subsection{Common Errors}
\label{subsec:mathematical-economics-common-errors}
%%%%%%%%%%%%%%%%%%%%%%%%%%%%%%%%%%%%%%%%%%%%%%%%%%%%%%%%%%%%

\subsubsection{Confusing Average and Marginal Quantities}

Average cost is

\[
AC(q)
=
\frac{C(q)}{q}.
\]

Marginal cost is

\[
MC(q)
=
C'(q).
\]

They measure different concepts.

%%%%%%%%%%%%%%%%%%%%%%%%%%%%%%%%%%%%%%%%%%%%%%%%%%%%%%%%%%%%
\subsubsection{Confusing Revenue and Profit}

Revenue is income from sales.

Profit subtracts cost:

\[
\boxed{
\pi=R-C.
}
\]

%%%%%%%%%%%%%%%%%%%%%%%%%%%%%%%%%%%%%%%%%%%%%%%%%%%%%%%%%%%%
\subsubsection{Applying \(MR=MC\) Without Checking a Maximum}

The condition

\[
MR=MC
\]

is typically a first-order condition.

One must also check:

\[
\boxed{
\text{second-order conditions},
}
\]

\[
\boxed{
\text{boundaries},
}
\]

and

\[
\boxed{
\text{constraints}.
}
\]

%%%%%%%%%%%%%%%%%%%%%%%%%%%%%%%%%%%%%%%%%%%%%%%%%%%%%%%%%%%%
\subsubsection{Confusing Utility With Measurable Happiness}

Utility functions usually represent preference rankings.

Their numerical levels often have no direct cardinal interpretation.

%%%%%%%%%%%%%%%%%%%%%%%%%%%%%%%%%%%%%%%%%%%%%%%%%%%%%%%%%%%%
\subsubsection{Treating Utility Transformations as Changing Preferences}

If \(g\) is strictly increasing, then

\[
\boxed{
u
}
\]

and

\[
\boxed{
g\circ u
}
\]

represent the same ordinal preferences.

%%%%%%%%%%%%%%%%%%%%%%%%%%%%%%%%%%%%%%%%%%%%%%%%%%%%%%%%%%%%
\subsubsection{Ignoring Corner Solutions}

The condition

\[
MRS
=
\frac{p_1}{p_2}
\]

applies to an interior optimum.

If a consumer chooses zero of one good, KKT or boundary conditions are
needed.

%%%%%%%%%%%%%%%%%%%%%%%%%%%%%%%%%%%%%%%%%%%%%%%%%%%%%%%%%%%%
\subsubsection{Confusing Marshallian and Hicksian Demand}

Marshallian demand holds income fixed.

Hicksian demand holds utility fixed.

They answer different questions.

%%%%%%%%%%%%%%%%%%%%%%%%%%%%%%%%%%%%%%%%%%%%%%%%%%%%%%%%%%%%
\subsubsection{Confusing Substitution and Income Effects}

The substitution effect changes consumption because relative prices change.

The income effect arises because purchasing power changes.

%%%%%%%%%%%%%%%%%%%%%%%%%%%%%%%%%%%%%%%%%%%%%%%%%%%%%%%%%%%%
\subsubsection{Confusing Returns to Scale With Marginal Returns}

Returns to scale change all inputs proportionally.

Diminishing marginal product changes one input while others remain fixed.

These are different concepts.

%%%%%%%%%%%%%%%%%%%%%%%%%%%%%%%%%%%%%%%%%%%%%%%%%%%%%%%%%%%%
\subsubsection{Assuming Cobb--Douglas Is Universal}

Cobb--Douglas functions are convenient models.

Real preferences and technologies may not have Cobb--Douglas structure.

%%%%%%%%%%%%%%%%%%%%%%%%%%%%%%%%%%%%%%%%%%%%%%%%%%%%%%%%%%%%
\subsubsection{Assuming Supply and Demand Curves Are Directly Observed Structural Laws}

Observed equilibrium data result from both supply and demand.

Structural relationships generally require identification assumptions.

%%%%%%%%%%%%%%%%%%%%%%%%%%%%%%%%%%%%%%%%%%%%%%%%%%%%%%%%%%%%
\subsubsection{Confusing Equilibrium With Social Optimality}

A market equilibrium need not maximize social welfare when there are:

\[
\boxed{
\text{externalities},
}
\]

\[
\boxed{
\text{market power},
}
\]

\[
\boxed{
\text{information problems},
}
\]

or

\[
\boxed{
\text{public goods}.
}
\]

%%%%%%%%%%%%%%%%%%%%%%%%%%%%%%%%%%%%%%%%%%%%%%%%%%%%%%%%%%%%
\subsubsection{Assuming a Tax Burden Falls on the Party Legally Charged}

Economic incidence depends on behavioral responses and elasticities.

Legal incidence and economic incidence can differ.

%%%%%%%%%%%%%%%%%%%%%%%%%%%%%%%%%%%%%%%%%%%%%%%%%%%%%%%%%%%%
\subsubsection{Ignoring Strategic Interaction}

In oligopoly or game-theoretic settings, an agent's optimum depends on others'
decisions.

Independent optimization can therefore be inappropriate.

%%%%%%%%%%%%%%%%%%%%%%%%%%%%%%%%%%%%%%%%%%%%%%%%%%%%%%%%%%%%
\subsubsection{Confusing Nash Equilibrium With Cooperation}

A Nash equilibrium only requires that no player wants to deviate
unilaterally.

It need not maximize total welfare.

%%%%%%%%%%%%%%%%%%%%%%%%%%%%%%%%%%%%%%%%%%%%%%%%%%%%%%%%%%%%
\subsubsection{Assuming Nash Equilibrium Is Unique}

Many games have multiple equilibria.

Equilibrium selection may require additional modeling assumptions.

%%%%%%%%%%%%%%%%%%%%%%%%%%%%%%%%%%%%%%%%%%%%%%%%%%%%%%%%%%%%
\subsubsection{Confusing Risk Aversion With Dislike of Low Wealth}

Risk aversion concerns curvature of utility:

\[
u''(W)<0.
\]

A person can prefer greater wealth while still having different attitudes
toward uncertainty.

%%%%%%%%%%%%%%%%%%%%%%%%%%%%%%%%%%%%%%%%%%%%%%%%%%%%%%%%%%%%
\subsubsection{Using Expected Value Instead of Expected Utility}

With nonlinear utility,

\[
u(\mathbb{E}[W])
\neq
\mathbb{E}[u(W)]
\]

in general.

%%%%%%%%%%%%%%%%%%%%%%%%%%%%%%%%%%%%%%%%%%%%%%%%%%%%%%%%%%%%
\subsubsection{Confusing Discounting With Inflation}

Time discounting reflects the relative valuation of outcomes at different
times.

Inflation concerns changes in the price level.

The concepts may interact but are not identical.

%%%%%%%%%%%%%%%%%%%%%%%%%%%%%%%%%%%%%%%%%%%%%%%%%%%%%%%%%%%%
\subsubsection{Assuming a Steady State Is Automatically Stable}

A steady state merely satisfies the equilibrium equation.

Stability must be analyzed separately.

%%%%%%%%%%%%%%%%%%%%%%%%%%%%%%%%%%%%%%%%%%%%%%%%%%%%%%%%%%%%
\subsubsection{Using Continuous-Time Stability Rules in Discrete Time}

For discrete

\[
x_{t+1}=F(x_t),
\]

local stability involves

\[
|F'(x^*)|<1.
\]

For continuous

\[
\dot{x}=f(x),
\]

the one-dimensional condition is

\[
f'(x^*)<0.
\]

%%%%%%%%%%%%%%%%%%%%%%%%%%%%%%%%%%%%%%%%%%%%%%%%%%%%%%%%%%%%
\subsubsection{Ignoring General-Equilibrium Feedback}

Changing one market can alter prices and decisions elsewhere.

Partial-equilibrium conclusions may not survive general-equilibrium effects.

%%%%%%%%%%%%%%%%%%%%%%%%%%%%%%%%%%%%%%%%%%%%%%%%%%%%%%%%%%%%
\subsubsection{Confusing Calibration With Causal Estimation}

Calibration selects parameters to match targets.

It does not automatically establish causal interpretation.

%%%%%%%%%%%%%%%%%%%%%%%%%%%%%%%%%%%%%%%%%%%%%%%%%%%%%%%%%%%%
\subsubsection{Ignoring Identification}

An estimator cannot reliably recover a structural parameter that is not
identified from observable data.

%%%%%%%%%%%%%%%%%%%%%%%%%%%%%%%%%%%%%%%%%%%%%%%%%%%%%%%%%%%%
\subsubsection{Assuming a Good In-Sample Fit Validates an Economic Model}

Many models can fit the same historical data.

Out-of-sample performance and structural plausibility remain important.

%%%%%%%%%%%%%%%%%%%%%%%%%%%%%%%%%%%%%%%%%%%%%%%%%%%%%%%%%%%%
\subsubsection{Treating Economic Models as Exact Descriptions}

Economic models simplify behavior, institutions, and information.

Their conclusions are conditional on assumptions.

%%%%%%%%%%%%%%%%%%%%%%%%%%%%%%%%%%%%%%%%%%%%%%%%%%%%%%%%%%%%
\subsection{Applications}
\label{subsec:mathematical-economics-applications}
%%%%%%%%%%%%%%%%%%%%%%%%%%%%%%%%%%%%%%%%%%%%%%%%%%%%%%%%%%%%

\begin{applicationbox}

\textbf{Consumer Choice.}
Optimization models explain how prices, income, and preferences influence
demand.

\medskip

\textbf{Firm Behavior.}
Production, cost, and profit models describe input choice, output, and supply.

\medskip

\textbf{Market Equilibrium.}
Supply and demand systems determine equilibrium prices and quantities.

\medskip

\textbf{Industrial Organization.}
Game theory models monopoly, oligopoly, pricing, entry, and strategic
competition.

\medskip

\textbf{Public Economics.}
Optimization and welfare analysis study taxation, subsidies, public goods,
and externalities.

\medskip

\textbf{Growth Theory.}
Differential and difference equations model capital accumulation, technology,
and long-run economic development.

\medskip

\textbf{Macroeconomics.}
Dynamic optimization describes consumption, saving, investment, and policy.

\medskip

\textbf{Mechanism Design.}
Game theory and optimization design auctions and institutions under strategic
behavior.

\medskip

\textbf{Financial Economics.}
Expected utility, portfolio choice, risk, and intertemporal decisions connect
directly with Mathematical Finance.

\medskip

\textbf{Economic Networks.}
Graph methods describe production chains, trade, financial exposure,
information diffusion, and network externalities.

\end{applicationbox}

%%%%%%%%%%%%%%%%%%%%%%%%%%%%%%%%%%%%%%%%%%%%%%%%%%%%%%%%%%%%
\subsection{Exercises}
\label{subsec:mathematical-economics-exercises}
%%%%%%%%%%%%%%%%%%%%%%%%%%%%%%%%%%%%%%%%%%%%%%%%%%%%%%%%%%%%

\begin{exercise}[1]
Define marginal cost.
\end{exercise}

\begin{exercise}[1]
Define marginal revenue.
\end{exercise}

\begin{exercise}[1]
Define elasticity.
\end{exercise}

\begin{exercise}[1]
Define a utility function.
\end{exercise}

\begin{exercise}[1]
Define marginal utility.
\end{exercise}

\begin{exercise}[1]
Define a budget constraint.
\end{exercise}

\begin{exercise}[1]
Define a production function.
\end{exercise}

\begin{exercise}[1]
Define marginal product.
\end{exercise}

\begin{exercise}[1]
Define market equilibrium.
\end{exercise}

\begin{exercise}[1]
Define Nash equilibrium.
\end{exercise}

\begin{exercise}[2]
Given

\[
C(q)=100+4q+q^2,
\]

compute marginal cost.
\end{exercise}

\begin{exercise}[2]
Given

\[
R(q)=20q-q^2,
\]

compute marginal revenue.
\end{exercise}

\begin{exercise}[2]
Find the profit-maximizing output for specified revenue and cost functions.
\end{exercise}

\begin{exercise}[2]
Compute price elasticity for

\[
Q=120-3P
\]

at a specified price.
\end{exercise}

\begin{exercise}[2]
Classify demand as elastic, inelastic, or unit elastic.
\end{exercise}

\begin{exercise}[2]
For

\[
u(x,y)=xy,
\]

compute marginal utilities.
\end{exercise}

\begin{exercise}[2]
Compute the marginal rate of substitution.
\end{exercise}

\begin{exercise}[2]
Write the Lagrangian for a two-good utility-maximization problem.
\end{exercise}

\begin{exercise}[3]
Solve

\[
\max_{x,y}
x^{1/2}y^{1/2}
\]

subject to

\[
p_xx+p_yy=m.
\]
\end{exercise}

\begin{exercise}[3]
Derive Marshallian demand for general Cobb--Douglas utility

\[
u(x,y)=x^\alpha y^\beta.
\]
\end{exercise}

\begin{exercise}[3]
Derive the indirect utility function for the previous problem.
\end{exercise}

\begin{exercise}[3]
Formulate the corresponding expenditure-minimization problem.
\end{exercise}

\begin{exercise}[3]
Derive Hicksian demand for a Cobb--Douglas consumer.
\end{exercise}

\begin{exercise}[3]
Explain income and substitution effects.
\end{exercise}

\begin{exercise}[2]
For

\[
q=AK^\alpha L^\beta,
\]

compute marginal products.
\end{exercise}

\begin{exercise}[2]
Determine returns to scale from \(\alpha+\beta\).
\end{exercise}

\begin{exercise}[3]
Derive the marginal rate of technical substitution.
\end{exercise}

\begin{exercise}[3]
Formulate a cost-minimization problem for a Cobb--Douglas firm.
\end{exercise}

\begin{exercise}[3]
Derive conditional factor demands.
\end{exercise}

\begin{exercise}[2]
Compute average and marginal cost for a specified cost function.
\end{exercise}

\begin{exercise}[3]
Show that average cost decreases when \(MC<AC\).
\end{exercise}

\begin{exercise}[2]
Solve a competitive market equilibrium from linear demand and supply.
\end{exercise}

\begin{exercise}[3]
Compute comparative statics of equilibrium price with respect to a demand
parameter.
\end{exercise}

\begin{exercise}[2]
Compute consumer surplus from a linear demand curve.
\end{exercise}

\begin{exercise}[2]
Compute producer surplus from a linear supply curve.
\end{exercise}

\begin{exercise}[3]
Analyze a per-unit tax in a linear market.
\end{exercise}

\begin{exercise}[3]
Compute buyer and seller tax incidence.
\end{exercise}

\begin{exercise}[3]
Explain the role of elasticity in tax incidence.
\end{exercise}

\begin{exercise}[2]
Analyze a binding price ceiling.
\end{exercise}

\begin{exercise}[2]
Analyze a binding price floor.
\end{exercise}

\begin{exercise}[2]
Derive marginal revenue for a monopolist.
\end{exercise}

\begin{exercise}[3]
Solve a monopoly problem with linear demand and constant marginal cost.
\end{exercise}

\begin{exercise}[3]
Compare monopoly and competitive outputs in the same market.
\end{exercise}

\begin{exercise}[2]
Define Cournot competition.
\end{exercise}

\begin{exercise}[3]
Derive reaction functions in a symmetric Cournot duopoly.
\end{exercise}

\begin{exercise}[3]
Solve the Cournot equilibrium.
\end{exercise}

\begin{exercise}[2]
Explain Bertrand competition.
\end{exercise}

\begin{exercise}[2]
Explain Stackelberg competition.
\end{exercise}

\begin{exercise}[3]
Solve a simple Stackelberg duopoly by backward induction.
\end{exercise}

\begin{exercise}[2]
Find Nash equilibria of a \(2\times2\) normal-form game.
\end{exercise}

\begin{exercise}[2]
Identify dominant strategies.
\end{exercise}

\begin{exercise}[3]
Compute a mixed-strategy Nash equilibrium of a simple game.
\end{exercise}

\begin{exercise}[3]
Explain why repeated interaction can change incentives.
\end{exercise}

\begin{exercise}[2]
Define a public good.
\end{exercise}

\begin{exercise}[2]
Define an externality.
\end{exercise}

\begin{exercise}[3]
Construct a model in which pollution creates a difference between private
and social marginal cost.
\end{exercise}

\begin{exercise}[3]
Compute a Pigouvian tax in a simple linear externality model.
\end{exercise}

\begin{exercise}[2]
Define a common-pool resource.
\end{exercise}

\begin{exercise}[2]
Define expected utility.
\end{exercise}

\begin{exercise}[2]
Determine whether a utility function is risk averse by inspecting
\(u''(W)\).
\end{exercise}

\begin{exercise}[3]
Compute the certainty equivalent of a simple lottery.
\end{exercise}

\begin{exercise}[3]
Compute the corresponding risk premium.
\end{exercise}

\begin{exercise}[3]
Compute absolute and relative Arrow--Pratt risk-aversion measures.
\end{exercise}

\begin{exercise}[2]
Write an intertemporal budget constraint.
\end{exercise}

\begin{exercise}[3]
Derive the consumption Euler equation.
\end{exercise}

\begin{exercise}[2]
Write a Bellman equation for a simple saving problem.
\end{exercise}

\begin{exercise}[3]
Explain dynamic programming in an economic setting.
\end{exercise}

\begin{exercise}[2]
Define a steady state.
\end{exercise}

\begin{exercise}[3]
Test local stability of a one-dimensional discrete economic model.
\end{exercise}

\begin{exercise}[3]
Test local stability of a one-dimensional continuous economic model.
\end{exercise}

\begin{exercise}[2]
Write the Solow capital-accumulation equation.
\end{exercise}

\begin{exercise}[3]
Solve for the Solow steady state with

\[
f(k)=k^\alpha.
\]
\end{exercise}

\begin{exercise}[3]
Analyze how the saving rate affects the Solow steady state.
\end{exercise}

\begin{exercise}[2]
Write the Leontief input-output system.
\end{exercise}

\begin{exercise}[3]
Compute production requirements for a small input-output matrix.
\end{exercise}

\begin{exercise}[3]
Explain the economic meaning of the Leontief inverse.
\end{exercise}

\begin{exercise}[2]
Define excess demand.
\end{exercise}

\begin{exercise}[2]
State Walras' law.
\end{exercise}

\begin{exercise}[3]
Explain why fixed-point mathematics appears in general equilibrium.
\end{exercise}

\begin{exercise}[2]
Explain a first-price auction.
\end{exercise}

\begin{exercise}[2]
Explain a second-price auction.
\end{exercise}

\begin{exercise}[3]
Explain why truthful bidding is weakly dominant in a standard second-price
private-value auction.
\end{exercise}

\begin{exercise}[2]
Define incentive compatibility.
\end{exercise}

\begin{exercise}[2]
Define mechanism design.
\end{exercise}

\begin{exercise}[2]
Distinguish structural and reduced-form economic models.
\end{exercise}

\begin{exercise}[3]
Explain identification.
\end{exercise}

\begin{exercise}[3]
Explain why supply and demand estimation can suffer from simultaneity.
\end{exercise}

\begin{exercise}[2]
Define calibration.
\end{exercise}

\begin{exercise}[3]
Construct a sensitivity analysis for a simple growth model.
\end{exercise}

\begin{exercise}[2]
Define a network effect.
\end{exercise}

\begin{exercise}[3]
Construct a simple coordination game with multiple equilibria.
\end{exercise}

\begin{exercise}[3]
Explain rational expectations conceptually.
\end{exercise}

\begin{exercise}[3]
Distinguish equilibrium from prediction.
\end{exercise}

\begin{exercise}[3]
Distinguish positive and normative economic analysis.
\end{exercise}

\begin{exercise}[4]
Derive the Slutsky equation for a two-good consumer.
\end{exercise}

\begin{exercise}[4]
Study duality between utility maximization and expenditure minimization.
\end{exercise}

\begin{exercise}[4]
Derive a firm's cost function from a homogeneous production function.
\end{exercise}

\begin{exercise}[4]
Study Shephard's lemma conceptually.
\end{exercise}

\begin{exercise}[4]
Study Hotelling's lemma conceptually.
\end{exercise}

\begin{exercise}[4]
Analyze welfare effects of a tax using surplus.
\end{exercise}

\begin{exercise}[4]
Compare Cournot, Bertrand, Stackelberg, monopoly, and competitive outcomes
in a common linear-demand model.
\end{exercise}

\begin{exercise}[4]
Analyze a repeated prisoner's dilemma with discounting.
\end{exercise}

\begin{exercise}[4]
Construct a public-goods contribution game.
\end{exercise}

\begin{exercise}[4]
Solve a simple mechanism-design problem.
\end{exercise}

\begin{exercise}[4]
Study auction equilibrium under symmetric private values.
\end{exercise}

\begin{exercise}[4]
Derive a dynamic consumption-saving model using Bellman equations.
\end{exercise}

\begin{exercise}[4]
Solve a finite-horizon consumption problem by backward induction.
\end{exercise}

\begin{exercise}[4]
Analyze stability of the Solow model.
\end{exercise}

\begin{exercise}[4]
Formulate the Ramsey growth model as an optimal-control problem.
\end{exercise}

\begin{exercise}[4]
Apply Pontryagin's Maximum Principle to an economic growth problem.
\end{exercise}

\begin{exercise}[4]
Study an overlapping-generations model.
\end{exercise}

\begin{exercise}[4]
Analyze an input-output economy using spectral-radius conditions.
\end{exercise}

\begin{exercise}[4]
Study equilibrium existence using a fixed-point theorem conceptually.
\end{exercise}

\begin{exercise}[4]
Construct and estimate a simultaneous-equations model.
\end{exercise}

\begin{exercise}[4]
Analyze an economic network with strategic complementarities.
\end{exercise}

\begin{exercise}[4]
Develop a dynamic policy optimization model under uncertainty.
\end{exercise}

%%%%%%%%%%%%%%%%%%%%%%%%%%%%%%%%%%%%%%%%%%%%%%%%%%%%%%%%%%%%
\subsection{Conceptual Summary}
\label{subsec:mathematical-economics-summary}
%%%%%%%%%%%%%%%%%%%%%%%%%%%%%%%%%%%%%%%%%%%%%%%%%%%%%%%%%%%%

Mathematical economics converts economic assumptions into optimization,
equilibrium, strategic, and dynamic models.

Marginal analysis uses derivatives such as

\[
\boxed{
MC=C'(q)
}
\]

and

\[
\boxed{
MR=R'(q).
}
\]

Consumer choice solves

\[
\boxed{
\max_{\mathbf{x}}
u(\mathbf{x})
}
\]

subject to

\[
\boxed{
\mathbf{p}^T\mathbf{x}\leq m.
}
\]

Firm behavior can be modeled through cost minimization and profit
maximization.

Competitive equilibrium satisfies

\[
\boxed{
Q_D(P^*)=Q_S(P^*).
}
\]

Strategic interaction is modeled using game theory and Nash equilibrium.

Risky choice uses expected utility:

\[
\boxed{
\mathbb{E}[u(W)].
}
\]

Dynamic economic problems use difference equations, differential equations,
optimal control, and Bellman equations.

Growth models study the evolution of capital:

\[
\boxed{
\dot{k}
=
sf(k)-(n+\delta)k.
}
\]

Input-output models use linear algebra:

\[
\boxed{
\mathbf{x}
=
(I-A)^{-1}\mathbf{d}.
}
\]

General equilibrium links many markets simultaneously.

Thus mathematical economics combines

\[
\boxed{
\text{optimization}
+
\text{equilibrium}
+
\text{game theory}
+
\text{probability}
+
\text{dynamics}
+
\text{statistics}.
}
\]

Its central task is to determine what economic behavior follows logically
from specified assumptions, constraints, incentives, and interactions.

%%%%%%%%%%%%%%%%%%%%%%%%%%%%%%%%%%%%%%%%%%%%%%%%%%%%%%%%%%%%
\subsection{Section Connections}
\label{subsec:mathematical-economics-connections}
%%%%%%%%%%%%%%%%%%%%%%%%%%%%%%%%%%%%%%%%%%%%%%%%%%%%%%%%%%%%

\chapterconnection{
Mathematical Economics completes the Applied Mathematics and Modeling chapter
by combining optimization, probability, game theory, dynamical systems,
control, networks, and statistics in economic applications. Consumer and
producer models use constrained optimization, strategic markets use Game
Theory, dynamic growth and policy models use differential equations and
optimal control, risky choice connects with probability and Mathematical
Finance, input-output systems use linear algebra, and general-equilibrium
existence connects with topology and fixed-point theory. The next chapter,
Mathematical Reference, collects notation, formulas, identities, transforms,
distributions, theorems, terminology, and indexing material for use
throughout the textbook.
}

%%%%%%%%%%%%%%%%%%%%%%%%%%%%%%%%%%%%%%%%%%%%%%%%%%%%%%%%%%%%
% SECTION SOURCE TRACKING
%%%%%%%%%%%%%%%%%%%%%%%%%%%%%%%%%%%%%%%%%%%%%%%%%%%%%%%%%%%%

%------------------------------------------------------------
% 13.9 Mathematical Economics
%------------------------------------------------------------
%
% Source basis:
%   - Standard microeconomic theory.
%   - Standard mathematical economics.
%   - Standard game theory, growth theory, equilibrium theory,
%     input-output analysis, and dynamic economic optimization.
%
% Used for:
%   - marginal analysis;
%   - cost, revenue, and profit;
%   - elasticity;
%   - consumer preferences;
%   - utility;
%   - indifference curves;
%   - marginal utility;
%   - marginal rate of substitution;
%   - budget constraints;
%   - utility maximization;
%   - Marshallian demand;
%   - indirect utility;
%   - expenditure minimization;
%   - Hicksian demand;
%   - income and substitution effects;
%   - production functions;
%   - Cobb--Douglas production;
%   - marginal products;
%   - returns to scale;
%   - isoquants;
%   - marginal rate of technical substitution;
%   - cost minimization;
%   - cost functions;
%   - average and marginal cost;
%   - competitive-firm optimization;
%   - market demand and supply;
%   - competitive equilibrium;
%   - comparative statics;
%   - consumer and producer surplus;
%   - taxes and subsidies;
%   - price controls;
%   - monopoly;
%   - marginal revenue;
%   - price discrimination;
%   - oligopoly;
%   - Cournot competition;
%   - reaction functions;
%   - Bertrand competition;
%   - Stackelberg competition;
%   - Nash equilibrium;
%   - dominant and mixed strategies;
%   - repeated games;
%   - public goods;
%   - externalities;
%   - Pigouvian taxes;
%   - common resources;
%   - expected utility;
%   - risk aversion;
%   - certainty equivalents;
%   - risk premiums;
%   - Arrow--Pratt measures;
%   - intertemporal choice;
%   - consumption-saving decisions;
%   - dynamic optimization;
%   - Bellman equations;
%   - economic steady states;
%   - local stability;
%   - Solow growth;
%   - optimal growth;
%   - input-output economics;
%   - Leontief models;
%   - general equilibrium;
%   - excess demand;
%   - Walras' law;
%   - fixed-point concepts;
%   - auctions;
%   - mechanism design;
%   - incentive compatibility;
%   - social choice;
%   - econometric connections;
%   - structural and reduced forms;
%   - identification;
%   - simultaneity;
%   - policy optimization;
%   - economic networks;
%   - network effects;
%   - multiple equilibria;
%   - expectations;
%   - calibration;
%   - sensitivity analysis;
%   - positive and normative analysis;
%   - applications and exercises.
%
% Notes:
%   - This completes Chapter 13 — Applied Mathematics and Modeling.
%   - Chapter 14 — Mathematical Reference follows.
%
%%%%%%%%%%%%%%%%%%%%%%%%%%%%%%%%%%%%%%%%%%%%%%%%%%%%%%%%%%%%